%=======================================================================
%  Reto I.1 - El pendulo que se vacia
%  Modelo lagrangiano de un pendulo fisico de masa variable
%  Compilar con: pdflatex archivo.tex  (dos veces, por las referencias)
%=======================================================================
\documentclass[11pt,a4paper]{article}

\usepackage[T1]{fontenc}
\usepackage[utf8]{inputenc}
\usepackage[spanish,es-noshorthands,es-nolists,es-noquoting,es-nodecimaldot]{babel}
\usepackage{mathptmx}
\usepackage{microtype}
\usepackage[margin=2.4cm]{geometry}
\usepackage{amsmath,amssymb,amsthm}
\usepackage{booktabs}
\usepackage{siunitx}
\sisetup{output-decimal-marker={.},per-mode=symbol,range-units=single}
\usepackage{enumitem}
\usepackage{graphicx}
\usepackage{tikz}
\usepackage{pgfplots}
\usepackage{caption}
\usepackage{float}
\usepackage[table]{xcolor}
\usepackage[colorlinks=true,linkcolor=azulosc,citecolor=azulosc,urlcolor=azulosc]{hyperref}

\pgfplotsset{compat=1.18}
\usetikzlibrary{patterns,decorations.pathmorphing,arrows.meta,calc}

\definecolor{azulosc}{RGB}{20,60,120}
\definecolor{agua}{RGB}{120,175,220}
\definecolor{rojo}{RGB}{190,40,40}
\definecolor{verde}{RGB}{20,120,70}
\definecolor{naranja}{RGB}{210,120,20}
\definecolor{gris}{RGB}{110,110,110}

\pgfplotsset{
  every axis/.append style={
    width=8.2cm, height=6.0cm, grid=major,
    grid style={gris!25,very thin},
    tick label style={font=\footnotesize},
    label style={font=\small},
    legend style={font=\footnotesize,fill=white,fill opacity=0.85,text opacity=1,draw=gris!50},
    every axis plot/.append style={line width=0.9pt}
  }
}

\newcommand{\dd}{\mathrm{d}}
\newcommand{\eps}{\varepsilon}
\newtheorem{resultado}{Resultado}

\title{\vspace{-1.2cm}\textbf{El péndulo que se vacía}\\[2mm]
\large Modelo lagrangiano de un péndulo físico de masa variable\\
\large con vaciado tipo Torricelli}
\author{Reto I.1 --- Mecánica Analítica}
\date{}

%=======================================================================
%  DATOS NUMERICOS EMBEBIDOS (generados por el modelo de la Sec. 10)
%=======================================================================
\begin{filecontents*}[overwrite]{curvas.dat}
# t h_mm mw Leq T Twkb Tacc Tjet gamma eps
0 250 1.25664 0.633993 1.5973 12 12 12 0.0158493 0.000429868
0.685315 246.455 1.23882 0.635459 1.59915 11.9422 12.0598 11.883 0.0156945 0.000427384
1.37542 242.911 1.22101 0.636928 1.601 11.8849 12.1211 11.7662 0.0155394 0.000424829
2.06553 239.393 1.20332 0.63839 1.60283 11.8284 12.1835 11.6503 0.0153851 0.000422235
2.75564 235.901 1.18577 0.639845 1.60466 11.7728 12.2471 11.5353 0.0152317 0.000419602
3.44574 232.434 1.16834 0.641292 1.60647 11.718 12.3118 11.4212 0.015079 0.00041693
4.13585 228.993 1.15104 0.642732 1.60827 11.6641 12.3778 11.3081 0.0149272 0.00041422
4.82596 225.577 1.13387 0.644163 1.61006 11.6109 12.4449 11.1958 0.0147763 0.000411473
5.51606 222.187 1.11684 0.645587 1.61184 11.5585 12.5133 11.0844 0.0146261 0.000408689
6.20138 218.846 1.10004 0.646993 1.6136 11.5073 12.5825 10.9746 0.0144777 0.000405889
6.89149 215.508 1.08326 0.6484 1.61535 11.4565 12.6534 10.8649 0.0143291 0.000403033
7.58159 212.195 1.06661 0.649798 1.61709 11.4064 12.7257 10.7559 0.0141813 0.000400142
8.2717 208.907 1.05008 0.651189 1.61882 11.3571 12.7993 10.6479 0.0140342 0.000397216
8.96181 205.645 1.03369 0.65257 1.62054 11.3085 12.8742 10.5406 0.0138879 0.000394255
9.65191 202.409 1.01742 0.653942 1.62224 11.2606 12.9506 10.434 0.0137423 0.00039126
10.342 199.199 1.00128 0.655305 1.62393 11.2135 13.0284 10.3283 0.0135975 0.000388231
11.0321 196.014 0.985275 0.656659 1.62561 11.167 13.1077 10.2233 0.0134534 0.000385168
11.7174 192.877 0.969505 0.657994 1.62726 11.1216 13.1879 10.1198 0.0133111 0.000382094
12.4076 189.743 0.953754 0.65933 1.62891 11.0765 13.2701 10.0163 0.0131684 0.000378964
13.0977 186.635 0.938132 0.660655 1.63054 11.0321 13.354 9.91352 0.0130265 0.000375802
13.7878 183.553 0.922639 0.66197 1.63217 10.9883 13.4394 9.81145 0.0128853 0.000372606
14.4779 180.497 0.907275 0.663276 1.63378 10.9452 13.5264 9.71009 0.0127447 0.000369378
15.168 177.466 0.89204 0.664571 1.63537 10.9028 13.6151 9.60941 0.0126049 0.000366116
15.8581 174.46 0.876933 0.665856 1.63695 10.861 13.7056 9.50941 0.0124657 0.000362822
16.5482 171.481 0.861956 0.667131 1.63852 10.8198 13.7977 9.41008 0.0123272 0.000359494
17.2335 168.547 0.847211 0.668386 1.64006 10.7795 13.891 9.31208 0.0121903 0.000356157
17.9236 165.619 0.832491 0.66964 1.64159 10.7395 13.9867 9.21405 0.0120531 0.000352763
18.6137 162.716 0.8179 0.670882 1.64312 10.7002 14.0844 9.11665 0.0119165 0.000349336
19.3038 159.839 0.803438 0.672113 1.64462 10.6614 14.1839 9.01988 0.0117805 0.000345875
19.9939 156.987 0.789105 0.673334 1.64612 10.6232 14.2854 8.92373 0.0116452 0.00034238
20.684 154.162 0.774901 0.674542 1.64759 10.5856 14.3888 8.82818 0.0115105 0.00033885
21.3742 151.361 0.760825 0.67574 1.64905 10.5486 14.4944 8.73324 0.0113764 0.000335286
22.0643 148.587 0.746879 0.676925 1.6505 10.5121 14.602 8.63888 0.0112429 0.000331687
22.7496 145.857 0.733158 0.678091 1.65192 10.4764 14.7111 8.54575 0.0111108 0.000328077
23.4397 143.134 0.719469 0.679253 1.65334 10.441 14.8231 8.45253 0.0109785 0.000324406
24.1298 140.436 0.705909 0.680403 1.65473 10.4062 14.9373 8.35988 0.0108467 0.000320698
24.8199 137.764 0.692478 0.68154 1.65612 10.3719 15.0539 8.26779 0.0107155 0.000316952
25.51 135.118 0.679176 0.682665 1.65748 10.3381 15.1729 8.17624 0.0105848 0.000313169
26.2001 132.497 0.666003 0.683777 1.65883 10.3048 15.2943 8.08522 0.0104546 0.000309346
26.8902 129.902 0.652959 0.684876 1.66017 10.2721 15.4183 7.99474 0.010325 0.000305483
27.5803 127.333 0.640044 0.685963 1.66148 10.2398 15.5448 7.90478 0.0101959 0.000301579
28.2656 124.807 0.627347 0.687028 1.66277 10.2083 15.673 7.81595 0.0100683 0.000297661
28.9557 122.289 0.614689 0.688088 1.66405 10.1771 15.8049 7.727 0.00994018 0.000293672
29.6459 119.796 0.60216 0.689135 1.66532 10.1463 15.9396 7.63855 0.0098126 0.000289639
30.336 117.329 0.58976 0.690167 1.66657 10.116 16.0771 7.5506 0.00968551 0.00028556
31.0261 114.888 0.57749 0.691186 1.6678 10.0862 16.2176 7.46313 0.0095589 0.000281434
31.7162 112.472 0.565348 0.69219 1.66901 10.0568 16.3612 7.37613 0.00943275 0.00027726
32.4063 110.083 0.553335 0.69318 1.6702 10.028 16.5078 7.28961 0.00930706 0.000273035
33.0964 107.718 0.541451 0.694155 1.67137 9.99953 16.6577 7.20355 0.00918183 0.000268758
33.7865 105.38 0.529696 0.695116 1.67253 9.97156 16.8108 7.11795 0.00905703 0.000264428
34.4718 103.083 0.518151 0.696055 1.67366 9.94423 16.9663 7.03339 0.00893353 0.000260072
35.1619 100.795 0.506653 0.696985 1.67478 9.91715 17.1264 6.94869 0.00880958 0.000255629
35.852 98.5337 0.495285 0.6979 1.67588 9.8905 17.29 6.86442 0.00868605 0.000251125
36.5421 96.2977 0.484045 0.698799 1.67696 9.86429 17.4574 6.78059 0.00856291 0.000246558
37.2322 94.0873 0.472934 0.699682 1.67801 9.8385 17.6286 6.69719 0.00844018 0.000241926
37.9223 91.9025 0.461952 0.700548 1.67905 9.81313 17.8037 6.61421 0.00831783 0.000237226
38.6125 89.7434 0.4511 0.701398 1.68007 9.78818 17.9829 6.53164 0.00819585 0.000232455
39.3026 87.61 0.440376 0.702231 1.68107 9.76364 18.1663 6.44949 0.00807424 0.00022761
39.9879 85.5168 0.429855 0.703042 1.68204 9.73967 18.3527 6.36831 0.00795382 0.000222722
40.678 83.4346 0.419388 0.70384 1.68299 9.71594 18.5449 6.28697 0.00783291 0.000217718
41.3681 81.378 0.40905 0.704621 1.68393 9.6926 18.7417 6.20602 0.00771234 0.00021263
42.0582 79.3471 0.398842 0.705383 1.68484 9.66965 18.9432 6.12547 0.00759209 0.000207452
42.7483 77.3418 0.388762 0.706127 1.68573 9.6471 19.1496 6.0453 0.00747216 0.000202181
43.4384 75.3622 0.378812 0.706852 1.68659 9.62492 19.3611 5.96552 0.00735254 0.000196811
44.1285 73.4082 0.36899 0.707558 1.68743 9.60313 19.5778 5.88612 0.00723321 0.000191339
44.8186 71.48 0.359297 0.708243 1.68825 9.58171 19.8 5.8071 0.00711418 0.000185757
45.5039 69.5905 0.3498 0.708904 1.68904 9.5608 20.0261 5.72899 0.00699624 0.000180102
46.194 67.7133 0.340364 0.70955 1.68981 9.54011 20.2596 5.65071 0.00687775 0.000174286
46.8842 65.8619 0.331058 0.710174 1.69055 9.51979 20.4991 5.57279 0.00675952 0.000168344
47.5743 64.0361 0.32188 0.710776 1.69127 9.49981 20.7448 5.49524 0.00664153 0.000162268
48.2644 62.236 0.312832 0.711356 1.69196 9.4802 20.9969 5.41804 0.00652379 0.00015605
48.9545 60.4615 0.303913 0.711914 1.69262 9.46092 21.2557 5.34121 0.00640627 0.000149683
49.6446 58.7127 0.295122 0.712448 1.69325 9.442 21.5213 5.26473 0.00628897 0.000143157
50.3347 56.9895 0.286461 0.712958 1.69386 9.4234 21.7941 5.1886 0.00617188 0.000136465
51.02 55.3038 0.277987 0.713441 1.69443 9.40527 22.0723 5.11335 0.00605579 0.000129644
51.7101 53.6318 0.269583 0.713901 1.69498 9.38735 22.36 5.03793 0.00593909 0.000122589
52.4002 51.9854 0.261307 0.714336 1.6955 9.36974 22.6557 4.96285 0.00582256 0.000115335
53.0903 50.3648 0.253161 0.714744 1.69598 9.35246 22.9597 4.88811 0.0057062 0.00010787
53.7804 48.7698 0.245144 0.715124 1.69643 9.33549 23.2722 4.81373 0.00558999 0.000100182
54.4705 47.2004 0.237255 0.715476 1.69685 9.31883 23.5935 4.73969 0.00547393 9.22576e-05
55.1606 45.6568 0.229496 0.715799 1.69723 9.30248 23.9241 4.666 0.00535802 8.40808e-05
55.8508 44.1388 0.221866 0.716091 1.69758 9.28643 24.2643 4.59265 0.00524222 7.56364e-05
56.5361 42.6567 0.214416 0.71635 1.69788 9.27079 24.6119 4.52015 0.00512735 6.69692e-05
57.2262 41.1898 0.207043 0.716579 1.69815 9.25532 24.9722 4.4475 0.00501179 5.79399e-05
57.9163 39.7486 0.199798 0.716774 1.69839 9.24014 25.3433 4.37519 0.00489633 4.85889e-05
58.6064 38.3331 0.192683 0.716934 1.69858 9.22525 25.7255 4.30323 0.00478095 3.88953e-05
59.2965 36.9432 0.185697 0.717058 1.69872 9.21062 26.1194 4.23163 0.00466566 2.88372e-05
59.9866 35.579 0.17884 0.717144 1.69883 9.19627 26.5253 4.16038 0.00455044 1.83909e-05
60.6767 34.2405 0.172112 0.717192 1.69888 9.18217 26.9437 4.08949 0.00443529 7.53098e-06
61.3668 32.9276 0.165512 0.717199 1.69889 9.16833 27.3752 4.01897 0.0043202 3.76945e-06
62.0521 31.6493 0.159087 0.717164 1.69885 9.15484 27.8172 3.94929 0.00420595 1.54558e-05
62.7422 30.3875 0.152744 0.717085 1.69876 9.14149 28.2763 3.8795 0.00409096 2.77219e-05
63.4323 29.1515 0.146531 0.716961 1.69861 9.12838 28.7501 3.81009 0.00397601 4.05205e-05
64.1225 27.9411 0.140447 0.716788 1.6984 9.1155 29.2393 3.74106 0.0038611 5.38861e-05
64.8126 26.7563 0.134492 0.716565 1.69814 9.10284 29.7444 3.67243 0.00374623 6.78555e-05
65.5027 25.5973 0.128666 0.716291 1.69781 9.09039 30.266 3.60419 0.00363139 8.24675e-05
66.1928 24.4639 0.122969 0.715961 1.69742 9.07814 30.8049 3.53637 0.00351659 9.77631e-05
66.8829 23.3561 0.117401 0.715575 1.69696 9.06609 31.3618 3.46897 0.00340183 0.000113785
67.573 22.2741 0.111962 0.715128 1.69644 9.05422 31.9373 3.40199 0.00328712 0.000130579
68.2583 21.2249 0.106688 0.714623 1.69584 9.0426 32.5281 3.33592 0.00317327 0.000148066
68.9484 20.194 0.101506 0.714048 1.69515 9.03107 33.1431 3.26985 0.00305868 0.000166538
69.6385 19.1887 0.096453 0.713405 1.69439 9.01969 33.779 3.20425 0.00294418 0.000185927
70.3286 18.2091 0.091529 0.712689 1.69354 9.00845 34.4367 3.13913 0.00282978 0.000206282
71.0187 17.2552 0.086734 0.711898 1.6926 8.99733 35.1169 3.07453 0.00271552 0.000227653
71.7088 16.3269 0.082068 0.711027 1.69156 8.98633 35.8205 3.01046 0.00260142 0.000250092
72.3989 15.4243 0.077531 0.710073 1.69043 8.97543 36.5482 2.94693 0.00248751 0.000273645
73.0891 14.5474 0.073123 0.709032 1.68919 8.96462 37.301 2.88399 0.00237386 0.000298358
73.7744 13.7019 0.0688732 0.707908 1.68785 8.95395 38.074 2.82207 0.00226129 0.000324086
74.4645 12.8761 0.0647224 0.706681 1.68639 8.94326 38.8788 2.76035 0.00214831 0.000351222
75.1546 12.076 0.0607005 0.705354 1.6848 8.93261 39.7108 2.6993 0.00203577 0.000379612
75.8447 11.3015 0.0568076 0.703923 1.68309 8.92198 40.5707 2.63894 0.00192376 0.000409267
76.5348 10.5527 0.0530437 0.702384 1.68125 8.91135 41.459 2.57933 0.00181239 0.000440179
77.2249 9.82957 0.0494088 0.700732 1.67927 8.90071 42.3761 2.52049 0.00170179 0.000472317
77.915 9.1321 0.0459029 0.698964 1.67715 8.89004 43.3224 2.46248 0.0015921 0.000505624
78.6051 8.46029 0.0425261 0.697076 1.67489 8.87932 44.298 2.40534 0.00148347 0.000540007
79.2904 7.81855 0.0393003 0.695079 1.67249 8.86861 45.2957 2.34952 0.00137684 0.000575083
79.9805 7.19789 0.0361805 0.692944 1.66991 8.85774 46.3291 2.2943 0.00127093 0.000611154
80.6706 6.60289 0.0331898 0.69068 1.66719 8.84678 47.3907 2.24013 0.00116672 0.000647731
81.3608 6.03356 0.030328 0.688289 1.6643 8.83571 48.4794 2.18709 0.0010645 0.000684489
82.0509 5.4899 0.0275952 0.685771 1.66125 8.82453 49.5936 2.13524 0.000964558 0.000721021
82.741 4.97189 0.0249915 0.683128 1.65805 8.81323 50.7315 2.08468 0.000867242 0.000756819
83.4311 4.47955 0.0225167 0.680365 1.65469 8.80181 51.8902 2.03548 0.00077292 0.000791269
84.1212 4.01288 0.0201709 0.677488 1.65119 8.79028 53.0665 1.98776 0.000681998 0.000823636
84.8065 3.57484 0.0179691 0.674528 1.64758 8.77875 54.2479 1.94192 0.000595501 0.000852863
85.4966 3.15932 0.0158805 0.671457 1.64382 8.76706 55.446 1.89743 0.000512676 0.000878376
86.1867 2.76946 0.0139208 0.668311 1.63997 8.75535 56.6467 1.85473 0.000434622 0.00089884
86.8768 2.40526 0.0120902 0.665112 1.63604 8.74365 57.8429 1.81394 0.000361823 0.000913026
87.5669 2.06673 0.0103885 0.661883 1.63206 8.73203 59.0266 1.77521 0.000294757 0.000919619
88.257 1.75386 0.00881585 0.658653 1.62807 8.72056 60.1887 1.73864 0.000233876 0.000917257
88.9471 1.46665 0.00737221 0.655458 1.62412 8.70935 61.3192 1.7044 0.000179581 0.000904587
89.6372 1.20511 0.00605756 0.652335 1.62025 8.69849 62.4068 1.6726 0.000132197 0.000880338
90.3226 0.970786 0.0048797 0.649346 1.61653 8.68818 63.4328 1.6436 9.21996e-05 0.000843714
91.0127 0.760395 0.00382216 0.646493 1.61297 8.67839 64.3991 1.61711 5.91104e-05 0.000793376
91.7028 0.575668 0.00289362 0.643846 1.60967 8.66936 65.2854 1.5935 3.31509e-05 0.000729076
92.3929 0.416605 0.00209408 0.641451 1.60667 8.66122 66.0791 1.57288 1.40444e-05 0.000650827
93.083 0.283206 0.00142355 0.639354 1.60404 8.65411 66.7678 1.55538 1.30298e-06 0.000559191
93.7731 0.175471 0.000882012 0.637599 1.60184 8.64817 67.3402 1.5411 -5.78055e-06 0.000455328
94.4632 0.0933998 0.000469479 0.636223 1.60011 8.64352 67.7864 1.53013 -8.12647e-06 0.000341005
95.1533 0.0369929 0.000185947 0.635256 1.59889 8.64026 68.0983 1.52255 -6.84683e-06 0.000218544
95.8434 0.00625 3.14159e-05 0.634722 1.59822 8.63846 68.2702 1.5184 -3.20847e-06 9.11789e-05
\end{filecontents*}
\begin{filecontents*}[overwrite]{env_completo.dat}
# t amp_deg
1.6038 11.8658
3.21181 11.7359
4.82399 11.6102
6.4402 11.4885
8.06053 11.3708
9.68479 11.2568
11.313 11.1466
12.945 11.0399
14.581 10.9367
16.2205 10.837
17.8638 10.7405
19.5107 10.6473
21.161 10.5572
22.8149 10.4702
24.472 10.3861
26.1325 10.305
27.7962 10.2267
29.4629 10.1511
31.1326 10.0782
32.8053 10.008
34.4808 9.94032
36.159 9.87515
37.8397 9.81243
39.523 9.75211
41.2085 9.69412
42.8964 9.6384
44.5863 9.58492
46.2782 9.5336
47.9719 9.4844
49.6673 9.43727
51.3642 9.39214
53.0623 9.34896
54.7617 9.30767
56.4619 9.2682
58.1629 9.23049
59.8643 9.19446
61.5659 9.16003
63.2674 9.12711
64.9685 9.09559
66.6688 9.06537
68.3679 9.0363
70.0654 9.00824
71.7607 8.98101
73.4533 8.95442
75.1426 8.92826
76.8279 8.90228
78.5083 8.87625
80.183 8.84994
81.8512 8.82316
83.512 8.79583
85.1644 8.76804
86.8079 8.74014
88.442 8.71283
90.0667 8.68723
91.6828 8.66486
93.2914 8.64741
94.8949 8.63652
\end{filecontents*}
\begin{filecontents*}[overwrite]{env_naive.dat}
# t amp_deg
1.60394 12.142
3.21224 12.2905
4.8249 12.4459
6.44188 12.6085
8.06312 12.7788
9.68858 12.9572
11.3182 13.144
12.9519 13.3397
14.5897 13.5449
16.2315 13.7601
17.8772 13.9859
19.5267 14.223
21.1801 14.472
22.8372 14.7337
24.498 15.0089
26.1624 15.2985
27.8304 15.6036
29.5019 15.9251
31.1767 16.2642
32.8549 16.6223
34.5363 17.0006
36.221 17.4009
37.9087 17.8247
39.5995 18.274
41.2932 18.7509
42.9899 19.2576
44.6893 19.7968
46.3915 20.3712
48.0963 20.9842
49.8037 21.6391
51.5136 22.3401
53.226 23.0915
54.9406 23.8985
56.6575 24.7666
58.3766 25.7022
60.0979 26.7124
61.8212 27.8053
63.5465 28.9899
65.2737 30.2764
67.0027 31.6762
68.7336 33.2017
70.4662 34.8668
72.2006 36.6863
73.9367 38.6758
75.6746 40.8506
77.4142 43.2249
79.1557 45.809
80.8993 48.6061
82.6449 51.6067
84.393 54.7806
86.1437 58.0665
87.8974 61.36
89.6543 64.5034
91.4145 67.2846
93.1776 69.4552
94.943 70.7755
\end{filecontents*}
\begin{filecontents*}[overwrite]{ciclos.dat}
# t T_num
2.0064 1.60695
3.61543 1.61111
5.2286 1.61522
6.84585 1.61928
8.46712 1.62328
10.0924 1.62722
11.7215 1.6311
13.3545 1.63492
14.9913 1.63867
16.6318 1.64235
18.276 1.64597
19.9237 1.6495
21.575 1.65296
23.2296 1.65634
24.8876 1.65964
26.5489 1.66285
28.2133 1.66597
29.8808 1.669
31.5512 1.67193
33.2246 1.67476
34.9007 1.67749
36.5795 1.68011
38.2609 1.68262
39.9447 1.68501
41.6308 1.68728
43.3192 1.68941
45.0096 1.69141
46.7019 1.69327
48.396 1.69497
50.0918 1.69651
51.789 1.69787
53.4874 1.69905
55.187 1.70003
56.8874 1.70079
58.5884 1.70131
60.2899 1.70157
61.9914 1.70155
63.6928 1.70121
65.3937 1.70052
67.0937 1.69944
68.7923 1.69793
70.4893 1.69592
72.1839 1.69336
73.8757 1.69019
75.5639 1.68633
77.248 1.6817
78.9269 1.67625
80.6 1.66992
82.2663 1.66269
83.925 1.65459
85.5751 1.64576
87.2163 1.63645
88.848 1.62704
90.4706 1.6181
92.0848 1.61033
93.6922 1.60445
\end{filecontents*}
\begin{filecontents*}[overwrite]{trazaA.dat}
# t theta
0 12
0.017155 11.973
0.03431 11.892
0.057184 11.7
0.074339 11.495
0.091494 11.237
0.11437 10.815
0.13152 10.442
0.14868 10.021
0.17155 9.3899
0.18871 8.8669
0.20586 8.3039
0.22873 7.4954
0.24589 6.8492
0.26304 6.172
0.28592 5.2265
0.30307 4.4894
0.32023 3.7322
0.3431 2.6971
0.36026 1.9063
0.37741 1.107
0.40029 0.034383
0.41744 -0.77046
0.4346 -1.5716
0.45747 -2.6277
0.47463 -3.406
0.49178 -4.1687
0.51465 -5.1553
0.53181 -5.8681
0.54896 -6.5541
0.57184 -7.4218
0.58899 -8.0335
0.60615 -8.6087
0.62902 -9.3143
0.64618 -9.7945
0.66333 -10.23
0.6862 -10.739
0.70336 -11.064
0.72052 -11.339
0.74339 -11.626
0.76054 -11.78
0.7777 -11.881
0.80057 -11.932
0.81773 -11.908
0.83488 -11.83
0.85776 -11.643
0.87491 -11.442
0.89207 -11.189
0.91494 -10.774
0.93209 -10.405
0.94925 -9.99
0.97212 -9.3669
0.98928 -8.85
1.0064 -8.2933
1.0293 -7.4934
1.0465 -6.8538
1.0636 -6.1833
1.0865 -5.2466
1.1037 -4.5162
1.1208 -3.7655
1.1437 -2.739
1.1608 -1.9544
1.1837 -0.89541
1.2009 -0.096222
1.218 0.70318
1.2409 1.7631
1.258 2.549
1.2752 3.3231
1.2981 4.3309
1.3152 5.0641
1.3324 5.7742
1.3553 6.6795
1.3724 7.3235
1.3896 7.9341
1.4124 8.6919
1.4296 9.2147
1.4467 9.6957
1.4696 10.269
1.4868 10.644
1.5039 10.972
1.5268 11.332
1.544 11.542
1.5611 11.701
1.584 11.83
1.6011 11.865
1.6183 11.847
1.6412 11.739
1.6583 11.598
1.6755 11.403
1.6984 11.065
1.7155 10.753
1.7327 10.393
1.7555 9.8404
1.7727 9.3742
1.7898 8.8659
1.8127 8.1266
1.8299 7.5292
1.847 6.898
1.8699 6.0088
1.8871 5.3101
1.9042 4.5876
1.9271 3.5928
1.9443 2.8276
1.9614 2.0498
1.9843 0.99928
2.0014 0.20589
2.0186 -0.58822
2.0415 -1.642
2.0586 -2.4238
2.0758 -3.1945
2.0986 -4.1988
2.1158 -4.93
2.133 -5.6389
2.1558 -6.5436
2.173 -7.1879
2.1901 -7.7996
2.213 -8.5599
2.2302 -9.0853
2.2473 -9.5697
2.2702 -10.148
2.2873 -10.529
2.3102 -10.962
2.3274 -11.23
2.3445 -11.448
2.3674 -11.658
2.3846 -11.754
2.4017 -11.797
2.4246 -11.774
2.4417 -11.694
2.4589 -11.562
2.4818 -11.305
2.4989 -11.054
2.5161 -10.752
2.539 -10.276
2.5561 -9.8653
2.5733 -9.4101
2.5961 -8.7379
2.6133 -8.1878
2.6305 -7.601
2.6533 -6.7662
2.6705 -6.1044
2.6876 -5.4152
2.7105 -4.4593
2.7277 -3.7186
2.7448 -2.9615
2.7677 -1.932
2.7849 -1.1496
2.802 -0.36212
2.8249 0.6894
2.842 1.4746
2.8592 2.253
2.8821 3.2742
2.8992 4.023
2.9164 4.7535
2.9392 5.6935
2.9564 6.3687
2.9735 7.0151
2.9964 7.8272
3.0136 8.3954
3.0307 8.9258
3.0536 9.57
3.0708 10.003
3.0879 10.392
3.1108 10.836
3.1279 11.114
3.1451 11.341
3.168 11.565
3.1851 11.672
3.2023 11.728
3.2252 11.72
3.2423 11.653
3.2595 11.534
3.2823 11.295
3.2995 11.057
3.3167 10.77
3.3395 10.312
3.3567 9.9141
3.3738 9.4724
3.3967 8.818
3.4139 8.2809
3.4367 7.5076
3.4539 6.8882
3.471 6.2381
3.4939 5.3284
3.5111 4.6181
3.5282 3.8871
3.5511 2.8862
3.5683 2.1202
3.5854 1.3449
3.6083 0.30256
3.6254 -0.48097
3.6426 -1.2621
3.6655 -2.294
3.6826 -3.056
3.6998 -3.8041
3.7227 -4.7741
3.7398 -5.4769
3.757 -6.1548
3.7798 -7.015
3.797 -7.6236
3.8142 -8.198
3.837 -8.906
3.8542 -9.3906
3.8713 -9.8331
3.8942 -10.354
3.9114 -10.691
3.9285 -10.98
3.9514 -11.289
3.9686 -11.462
3.9857 -11.584
4.0086 -11.666
4.0257 -11.666
4.0429 -11.615
4.0658 -11.466
4.0829 -11.295
4.1001 -11.073
4.1229 -10.701
4.1401 -10.366
4.1573 -9.9844
4.1801 -9.4072
4.1973 -8.9252
4.2144 -8.4033
4.2373 -7.6497
4.2545 -7.0444
4.2716 -6.4077
4.2945 -5.5149
4.3117 -4.8163
4.3288 -4.0963
4.3517 -3.1086
4.3688 -2.3513
4.386 -1.5837
4.4089 -0.55004
4.426 0.22829
4.4432 1.0054
4.4661 2.0336
4.4832 2.7944
4.5004 3.5424
4.5232 4.5142
4.5404 5.2197
4.5575 5.9016
4.5804 6.769
4.5976 7.3844
4.6204 8.1528
4.6376 8.6868
4.6547 9.1819
4.6776 9.7777
4.6948 10.174
4.7119 10.524
4.7348 10.919
4.752 11.158
4.7691 11.347
4.792 11.521
4.8091 11.591
4.8263 11.61
4.8492 11.555
4.8663 11.453
4.8835 11.301
4.9064 11.02
4.9235 10.751
4.9407 10.435
4.9635 9.9422
4.9807 9.5205
4.9979 9.0565
5.0207 8.3756
5.0379 7.8212
5.055 7.232
5.0779 6.3969
5.0951 5.7371
5.1122 5.0518
5.1351 4.1037
5.1523 3.3711
5.1694 2.6235
5.1923 1.6094
5.2094 0.84024
5.2266 0.067483
5.2495 -0.96243
5.2666 -1.7301
5.2838 -2.4898
5.3067 -3.4845
5.3238 -4.2126
5.341 -4.9217
5.3638 -5.8322
5.381 -6.4848
5.3981 -7.1083
5.421 -7.8897
5.4382 -8.4349
5.4553 -8.9424
5.4782 -9.5565
5.4954 -9.9676
5.5125 -10.334
5.5354 -10.751
5.5525 -11.008
5.5697 -11.217
5.5926 -11.416
5.6097 -11.507
5.6269 -11.547
5.6498 -11.52
5.6669 -11.44
5.6841 -11.31
5.7069 -11.058
5.7241 -10.812
5.747 -10.409
5.7641 -10.053
5.7813 -9.6523
5.8041 -9.0518
5.8213 -8.5544
5.8385 -8.019
5.8613 -7.2503
5.8785 -6.6359
5.8956 -5.9922
5.9185 -5.093
5.9357 -4.392
5.9528 -3.6716
5.9757 -2.6863
5.9928 -1.9332
6.01 -1.1717
6.0329 -0.14901
6.05 0.61897
6.0672 1.384
6.0901 2.3936
6.1072 3.1385
6.1244 3.8693
6.1472 4.8159
6.1644 5.5011
6.1816 6.1616
6.2044 6.9988
6.2216 7.5906
6.2387 8.1485
6.2616 8.8354
6.2788 9.3049
6.2959 9.733
6.3188 10.236
6.336 10.561
6.3531 10.838
6.376 11.133
6.3931 11.297
6.4103 11.411
6.4332 11.484
6.4503 11.48
6.4675 11.424
6.4904 11.272
6.5075 11.1
6.5247 10.878
6.5475 10.508
6.5647 10.176
6.5818 9.7986
6.6047 9.2288
6.6219 8.7535
6.639 8.2396
6.6619 7.498
6.6791 6.9028
6.6962 6.2771
6.7191 5.4002
6.7362 4.7143
6.7534 4.0076
6.7763 3.0384
6.7934 2.2955
6.8106 1.5426
6.8335 0.52894
6.8506 -0.23429
6.8735 -1.2494
6.8906 -2.0045
6.9078 -2.7505
6.9307 -3.7252
6.9478 -4.4372
6.965 -5.1292
6.9878 -6.0158
7.005 -6.6497
7.0222 -7.2539
7.045 -8.0088
7.0622 -8.5337
7.0793 -9.0206
7.1022 -9.6073
7.1194 -9.9977
7.1365 -10.344
7.1594 -10.734
7.1765 -10.971
7.1937 -11.16
7.2166 -11.334
7.2337 -11.407
7.2509 -11.429
7.2738 -11.38
7.2909 -11.285
7.3081 -11.14
7.3309 -10.87
7.3481 -10.611
7.3653 -10.306
7.3881 -9.8286
7.4053 -9.4196
7.4224 -8.9691
7.4453 -8.3072
7.4625 -7.7677
7.4796 -7.1939
7.5025 -6.38
7.5197 -5.7364
7.5368 -5.0675
7.5597 -4.1415
7.5768 -3.4253
7.594 -2.6941
7.6169 -1.7014
7.634 -0.94788
7.6512 -0.19031
7.6741 0.82024
7.6912 1.5741
7.7084 2.3208
7.7312 3.2996
7.7484 4.0168
7.7656 4.716
7.7884 5.615
7.8056 6.2603
7.8227 6.8779
7.8456 7.6532
7.8628 8.1953
7.8799 8.7011
7.9028 9.3151
7.9199 9.7277
7.9371 10.097
7.96 10.52
7.9771 10.784
8 11.06
\end{filecontents*}
\begin{filecontents*}[overwrite]{trazaB.dat}
# t theta
88 -1.1547
88.017 -0.59262
88.034 -0.028075
88.056 0.72442
88.073 1.2855
88.09 1.8413
88.112 2.5697
88.129 3.1038
88.146 3.6248
88.168 4.2954
88.185 4.7775
88.202 5.2395
88.224 5.821
88.241 6.2288
88.258 6.6104
88.28 7.0757
88.297 7.3902
88.314 7.6736
88.336 8.0012
88.353 8.2078
88.37 8.3799
88.392 8.5546
88.409 8.6437
88.426 8.6965
88.448 8.7101
88.465 8.6776
88.482 8.6086
88.505 8.4604
88.521 8.3077
88.538 8.1201
88.561 7.817
88.578 7.5511
88.594 7.2534
88.617 6.8093
88.633 6.4426
88.65 6.0487
88.673 5.4841
88.69 5.0335
88.706 4.5616
88.729 3.9029
88.746 3.3893
88.763 2.8614
88.785 2.1392
88.802 1.5866
88.819 1.0274
88.841 0.27539
88.858 -0.29027
88.875 -0.8547
88.897 -1.6012
88.914 -2.1536
88.931 -2.6967
88.953 -3.4029
88.97 -3.916
88.987 -4.4125
89.009 -5.0451
89.026 -5.495
89.043 -5.9216
89.065 -6.4511
89.082 -6.8167
89.099 -7.1535
89.121 -7.5552
89.138 -7.8195
89.16 -8.1202
89.177 -8.3059
89.194 -8.4565
89.217 -8.6016
89.233 -8.6681
89.25 -8.6979
89.273 -8.6806
89.29 -8.6248
89.306 -8.5325
89.329 -8.3535
89.346 -8.1779
89.362 -7.9677
89.385 -7.6352
89.402 -7.348
89.418 -7.0297
89.441 -6.5592
89.458 -6.1737
89.475 -5.7621
89.497 -5.1755
89.514 -4.7097
89.531 -4.2239
89.553 -3.5486
89.57 -3.0244
89.587 -2.4872
89.609 -1.755
89.626 -1.1968
89.643 -0.63351
89.665 0.12129
89.682 0.68709
89.699 1.25
89.721 1.9918
89.738 2.5386
89.755 3.0746
89.777 3.7685
89.794 4.2705
89.811 4.7544
89.833 5.3677
89.85 5.8014
89.867 6.2103
89.889 6.7142
89.906 7.0591
89.923 7.374
89.945 7.7448
89.962 7.9847
89.979 8.1907
90.001 8.4111
90.018 8.5348
90.035 8.6222
90.058 8.6817
90.074 8.6834
90.091 8.6482
90.114 8.5441
90.13 8.4236
90.147 8.2673
90.17 8.0043
90.186 7.7672
90.203 7.4971
90.226 7.0875
90.243 6.7449
90.265 6.2438
90.282 5.8366
90.299 5.4046
90.321 4.793
90.338 4.3102
90.355 3.809
90.377 3.1158
90.394 2.5801
90.411 2.0334
90.433 1.2913
90.45 0.72799
90.467 0.16157
90.489 -0.59431
90.506 -1.1586
90.523 -1.7179
90.545 -2.4518
90.562 -2.9903
90.579 -3.516
90.601 -4.1931
90.618 -4.6803
90.635 -5.1476
90.657 -5.7359
90.674 -6.1489
90.691 -6.5355
90.713 -7.0074
90.73 -7.3265
90.747 -7.6143
90.77 -7.9473
90.786 -8.1576
90.803 -8.333
90.826 -8.5115
90.843 -8.603
90.859 -8.6577
90.882 -8.6732
90.898 -8.6416
90.915 -8.5731
90.938 -8.4248
90.955 -8.2715
90.971 -8.0829
90.994 -7.7778
91.011 -7.51
91.028 -7.2101
91.05 -6.7625
91.067 -6.3928
91.084 -5.9957
91.106 -5.4265
91.123 -4.9723
91.14 -4.4968
91.162 -3.8331
91.179 -3.3158
91.196 -2.7844
91.218 -2.0575
91.235 -1.5018
91.252 -0.93958
91.274 -0.18413
91.291 0.3837
91.308 0.94988
91.33 1.698
91.347 2.2509
91.37 2.9726
91.386 3.4994
91.403 4.011
91.425 4.666
91.442 5.1342
91.459 5.5803
91.482 6.1374
91.498 6.5248
91.515 6.8841
91.538 7.317
91.555 7.6052
91.571 7.8607
91.594 8.1488
91.611 8.3242
91.627 8.4638
91.65 8.5935
91.667 8.6477
91.683 8.6648
91.706 8.6298
91.723 8.5603
91.74 8.4541
91.762 8.2561
91.779 8.0661
91.796 7.8415
91.818 7.4897
91.835 7.1882
91.852 6.8557
91.874 6.3669
91.891 5.968
91.908 5.5435
91.93 4.9407
91.947 4.4634
91.964 3.967
91.986 3.2788
92.003 2.7459
92.02 2.2011
92.042 1.4604
92.059 0.89717
92.076 0.33009
92.098 -0.42781
92.115 -0.99439
92.132 -1.5567
92.154 -2.2955
92.171 -2.8384
92.188 -3.369
92.21 -4.0535
92.227 -4.5468
92.244 -5.0204
92.266 -5.6178
92.283 -6.0378
92.3 -6.4318
92.323 -6.9137
92.339 -7.2405
92.356 -7.5361
92.379 -7.8795
92.395 -8.0975
92.412 -8.2807
92.435 -8.4693
92.451 -8.5681
92.468 -8.6301
92.491 -8.6549
92.508 -8.6299
92.53 -8.5388
92.547 -8.4275
92.564 -8.2798
92.586 -8.0275
92.603 -7.7977
92.62 -7.5343
92.642 -7.1327
92.659 -6.7954
92.676 -6.4287
92.698 -5.8968
92.715 -5.468
92.732 -5.0155
92.754 -4.3788
92.771 -3.8788
92.788 -3.3621
92.81 -2.6509
92.827 -2.1038
92.844 -1.5475
92.866 -0.79589
92.883 -0.22781
92.9 0.34126
92.922 1.0973
92.939 1.6591
92.956 2.2137
92.978 2.9378
92.995 3.4663
93.012 3.9799
93.035 4.6373
93.051 5.1073
93.068 5.5551
93.091 6.1144
93.108 6.5033
93.124 6.8639
93.147 7.2983
93.163 7.5875
93.18 7.8438
93.203 8.1325
93.22 8.3082
93.236 8.4479
93.259 8.5772
93.276 8.6309
93.293 8.6473
93.315 8.611
93.332 8.5403
93.349 8.4326
93.371 8.2324
93.388 8.0405
93.405 7.8139
93.427 7.4592
93.444 7.1553
93.461 6.8204
93.483 6.3282
93.5 5.9267
93.517 5.4996
93.539 4.8933
93.556 4.4135
93.573 3.9146
93.595 3.2232
93.612 2.6881
93.635 1.9568
93.651 1.3981
93.668 0.83334
93.691 0.075092
93.707 -0.49429
93.724 -1.0615
93.747 -1.8102
93.763 -2.3629
93.78 -2.9053
93.803 -3.6084
93.82 -4.1177
93.836 -4.6092
93.859 -5.2329
93.876 -5.6744
93.892 -6.0912
93.915 -6.6056
93.932 -6.9581
93.948 -7.2804
93.971 -7.6607
93.988 -7.9073
94.004 -8.1195
94.027 -8.3476
94.044 -8.4765
94.061 -8.5686
94.083 -8.6335
94.1 -8.6385
94.117 -8.6061
94.139 -8.5047
94.156 -8.3856
94.173 -8.2301
94.195 -7.9673
94.212 -7.7297
94.229 -7.4585
94.251 -7.0468
94.268 -6.702
94.285 -6.3281
94.307 -5.7871
94.324 -5.3518
94.341 -4.8931
94.363 -4.2488
94.38 -3.7437
94.397 -3.2223
94.419 -2.5056
94.436 -1.9551
94.453 -1.396
94.475 -0.64154
94.492 -0.072108
94.509 0.49764
94.531 1.2535
94.548 1.8143
94.565 2.3672
94.588 3.088
94.604 3.6132
94.621 4.1226
94.644 4.7735
94.66 5.2377
94.677 5.6791
94.7 6.2288
94.716 6.6096
94.739 7.0724
94.756 7.3838
94.773 7.663
94.795 7.9833
94.812 8.1831
94.829 8.3473
94.851 8.5097
94.868 8.5884
94.885 8.6297
94.907 8.6265
94.924 8.5803
94.941 8.4967
94.963 8.328
94.98 8.159
94.997 7.9546
95.019 7.6284
95.036 7.3448
95.053 7.0293
95.075 6.5612
95.092 6.1766
95.109 5.7651
95.131 5.1776
95.148 4.7104
95.165 4.2226
95.188 3.544
95.204 3.0167
95.221 2.4762
95.243 1.7392
95.26 1.1773
95.277 0.61019
95.3 -0.14964
95.316 -0.71908
95.333 -1.2854
95.356 -2.0313
95.373 -2.5807
95.389 -3.1188
95.412 -3.8147
95.428 -4.3175
95.445 -4.8014
95.468 -5.4138
95.485 -5.8457
95.501 -6.2522
95.524 -6.7514
95.541 -7.0917
95.558 -7.4011
95.58 -7.7633
95.597 -7.9955
95.614 -8.193
95.636 -8.4007
95.653 -8.5139
95.67 -8.5899
95.692 -8.6332
95.709 -8.6218
95.726 -8.5729
95.748 -8.4496
95.765 -8.3142
95.782 -8.1425
95.804 -7.8587
95.821 -7.6057
95.843 -7.2171
\end{filecontents*}
\begin{filecontents*}[overwrite]{subito.dat}
# tv_over_T a0 a90 a45
0.0125211 11.994 6.22847 9.41459
0.0162082 11.99 6.23343 9.37068
0.0209809 11.9832 6.24171 9.31339
0.027159 11.9719 6.25553 9.23855
0.0351563 11.953 6.27849 9.14073
0.0455086 11.9215 6.31646 9.01301
0.0589093 11.8692 6.37868 8.84687
0.0762559 11.7828 6.47917 8.63279
0.0987105 11.6419 6.63765 8.3625
0.127777 11.4159 6.87817 8.03544
0.165403 11.0654 7.2215 7.67359
0.214108 10.5539 7.66527 7.348
0.277156 9.89126 8.14905 7.19914
0.358768 9.22769 8.52212 7.36835
0.464412 8.88831 8.59788 7.76361
0.601165 8.94725 8.46065 7.9643
0.778186 8.78218 8.63004 7.9307
1.00733 8.62451 8.6881 8.29278
1.30396 8.72667 8.60546 8.19683
1.68793 8.71514 8.57958 8.42308
2.18496 8.66014 8.62855 8.4831
2.82835 8.67517 8.61362 8.48319
3.6612 8.62123 8.66482 8.52327
4.73929 8.62862 8.65592 8.55233
6.13484 8.63438 8.64968 8.56335
7.94133 8.6344 8.64883 8.58913
10.2798 8.63419 8.64883 8.59764
13.3068 8.63477 8.64802 8.60951
17.2252 8.633 8.64967 8.61816
22.2974 8.63333 8.64927 8.62373
28.8631 8.63361 8.64894 8.62899
37.3623 8.63339 8.64914 8.63246
48.3641 8.63331 8.64921 8.6354
62.6056 8.63334 8.64917 8.63773
\end{filecontents*}

\begin{document}
\maketitle

\begin{abstract}
\noindent
Se construye un modelo lagrangiano completo de una botella cilíndrica parcialmente llena de agua
que oscila como péndulo mientras se vacía por un orificio de área $a$ practicado en su fondo.
La botella se trata como \emph{sólido rígido} (péndulo físico), de modo que el lagrangiano contiene
la energía cinética de traslación del centro de masa, la energía cinética de rotación en torno a él
y la energía potencial gravitatoria del centro de masa. El gasto se modela con la ley de Torricelli
para un tanque cilíndrico, lo que produce la ley de vaciado $h(t)=h_0(1-t/t_v)^2$.
Se muestra que el lagrangiano ingenuo con masa dependiente del tiempo es \emph{incompleto}, y se
corrige mediante el balance de momento angular con flujo de masa. La ecuación resultante,
$I_O\ddot\theta+\Gamma\dot\theta+g\,S\sin\theta=0$, se resuelve en la aproximación adiabática (WKB)
y se integra numéricamente. Las conclusiones centrales son: \textbf{(i)} el período
\emph{no} es monótono: crece, alcanza un máximo cuando la superficie libre cruza el centro de
oscilación instantáneo, y decrece de nuevo hasta el valor de la botella vacía; \textbf{(ii)} la
amplitud obedece $\Theta\propto (S\,I_O)^{-1/4}\exp(-\tfrac12\!\int\!\dot m_{\text{out}}s^2/I_O\,\dd t)$,
producto de un factor de bombeo adiabático y un factor de frenado por el chorro; y \textbf{(iii)} el
factor total de cambio de amplitud a lo largo de un vaciado completo es \emph{independiente del
tamaño del orificio}, resultado que constituye la predicción experimental más limpia del modelo.
\end{abstract}

\tableofcontents
\bigskip

%=======================================================================
\section{Planteamiento e idealizaciones}
%=======================================================================

Una botella cilíndrica de altura $H$ y radio interior $R$ cuelga de una cuerda de longitud $L_c$
y oscila en un plano vertical. El agua escapa por un orificio circular de área $a$ situado en el
centro del fondo. Se pregunta cómo evolucionan el \textbf{período} $T$ y la \textbf{amplitud}
$\Theta$ mientras la masa disminuye.

El problema es interesante precisamente porque tres efectos compiten y ninguno domina de forma obvia:

\begin{enumerate}[label=(\alph*),leftmargin=*]
\item \textbf{La distribución de masa cambia de forma.} El agua remanente se acumula siempre en el
fondo, de modo que el centro de masa del agua \emph{baja} (se aleja del pivote) aunque su masa
disminuya. El centro de masa del conjunto describe una trayectoria no monótona.
\item \textbf{El sistema pierde momento angular.} El chorro sale por el punto más alejado del
pivote y se lleva consigo momento angular que el resto del sistema no recupera.
\item \textbf{La variación es lenta.} Si $t_v \gg T$ el sistema es un oscilador de parámetros
lentamente variables y aparece un invariante adiabático que fija la amplitud.
\end{enumerate}

\subsection{Elección de ruta de modelado}

De entre las rutas propuestas en el enunciado, este trabajo adopta las siguientes decisiones,
que se justifican cuantitativamente en la Sección~\ref{sec:validez}:

\begin{center}
\renewcommand{\arraystretch}{1.25}
\begin{tabular}{@{}p{6.6cm}p{9.2cm}@{}}
\toprule
\textbf{Decisión} & \textbf{Justificación} \\
\midrule
Péndulo \textbf{físico} (sólido rígido), no masa puntual & Para una masa puntual el período y la
amplitud son \emph{exactamente constantes} durante el vaciado (Sec.~\ref{sec:puntual}): todo el
fenómeno es un efecto de cuerpo extenso. \\
Se \textbf{incluye} el desplazamiento del centro de masa & Es el origen físico de la variación del
período; ignorarlo elimina el efecto. \\
Gasto tipo \textbf{Torricelli} con coeficiente de descarga $C_d$ & El gasto constante es una mala
aproximación: $|\dot m|$ decrece linealmente hasta anularse, y con él el frenado
$\Gamma\propto|\dot h|$. \\
El agua sale con velocidad relativa \textbf{puramente axial} & El orificio está sobre el eje que
pasa por el pivote, luego el chorro no ejerce torque, pero sí extrae momento angular
(Sec.~\ref{sec:masavar}). \\
Movimiento interno del líquido \textbf{despreciado} (agua ``congelada'') & El primer modo de
chapoteo está fuera de resonancia por un factor $\gtrsim 5$ salvo en el último milímetro de columna
(Sec.~\ref{sec:validez}). \\
Vaciado lento tratado como proceso \textbf{adiabático} & El parámetro de adiabaticidad es
$\eps=|\dot\omega|/\omega^2\lesssim 10^{-3}$ para un orificio de \SI{5}{\milli\meter}. \\
\bottomrule
\end{tabular}
\end{center}

\subsection{Idealizaciones explícitas}
\label{sec:ideal}

\begin{enumerate}[label=\textbf{I\arabic*.},leftmargin=*,itemsep=2pt]
\item La cuerda es inextensible, de masa despreciable y permanece \textbf{tensa y recta}.
\item La cuerda y el eje de la botella permanecen \textbf{colineales}: el conjunto
cuerda\,+\,botella se comporta como un único sólido rígido articulado en el pivote $O$.
\item El movimiento ocurre en un \textbf{plano vertical fijo} (no hay componente azimutal ni
torsión respecto del eje vertical).
\item La botella es un \textbf{cilindro recto} de sección interior constante $A=\pi R^2$.
\item El agua se mueve como un \textbf{sólido rígido solidario a la botella} (hipótesis de agua
congelada): se desprecia el campo de velocidades interno frente a la velocidad de arrastre.
\item La \textbf{superficie libre} permanece plana y perpendicular al eje de la botella.
\item El fluido es \textbf{incompresible} y el flujo por el orificio es \textbf{cuasiestacionario}.
\item Se desprecian la resistencia del aire y la disipación en el pivote (se discuten en la
Sec.~\ref{sec:experimento}).
\end{enumerate}

%=======================================================================
\section{Vínculos y grados de libertad}
%=======================================================================

Consideremos el sistema material formado por la botella (rígida, masa $m_b$) más el agua contenida
en su interior en el instante $t$. Antes de imponer vínculos, la botella como cuerpo rígido en el
espacio posee $6$ grados de libertad y el agua, como medio continuo, posee infinitos.

\subsection{Enumeración de los vínculos}

\begin{description}[leftmargin=1.2cm,style=nextline,itemsep=3pt]

\item[V1 --- Cuerda inextensible y tensa (holónomo, escleronómico).]
Todo punto material $P$ del conjunto rígido mantiene su distancia al pivote:
\begin{equation}
|\mathbf r_P - \mathbf r_O| = \text{cte}, \qquad \forall P .
\end{equation}
Esto elimina los tres grados de libertad de traslación del cuerpo. Es un vínculo \emph{unilateral}
en rigor ($|\mathbf r|\le L_c$ para la cuerda), que se comporta como bilateral mientras la tensión
sea positiva; esto exige $\dot\theta^2$ moderado, condición satisfecha para amplitudes
$\Theta \lesssim \ang{60}$.

\item[V2 --- Alineación cuerda--botella (holónomo).]
El eje de simetría de la botella contiene al pivote. Elimina dos grados de libertad de rotación
(cabeceo y balanceo relativos a la cuerda) y hace que el sistema sea un péndulo \emph{compuesto} de
un solo cuerpo efectivo.

\item[V3 --- Movimiento plano (holónomo).]
$z\equiv 0$ para todo punto: el movimiento se restringe al plano vertical $xy$. Elimina el grado de
libertad azimutal (péndulo esférico $\to$ péndulo plano) y la rotación propia en torno al eje de
la botella.

\item[V4 --- Rigidez del agua respecto de la botella (holónomo, \emph{reónomo}).]
Cada partícula de fluido tiene velocidad igual a la de arrastre del sólido más una deriva axial
uniforme $\mathbf v_{\text{der}} = -(Q/A)\,\hat{\mathbf e}_r$. En la aproximación de agua congelada
se desprecia $\mathbf v_{\text{der}}$ en la cinemática rotacional. Este vínculo es
\textbf{reónomo}: depende explícitamente de $t$ a través de $h(t)$.

\item[V5 --- Superficie libre plana y horizontal en el marco de la botella (holónomo, reónomo).]
Fija la geometría de la región ocupada por el agua a partir de un único parámetro $h(t)$.

\item[V6 --- Ligadura de vaciado (reónoma, impuesta externamente).]
La altura de la columna sigue una ley prescrita
$h = h(t)$ determinada por la hidrodinámica del orificio (Sec.~\ref{sec:torricelli}). Éste no es
un vínculo mecánico ordinario sino una \emph{ligadura constitutiva}: convierte el sistema en uno de
masa variable y hace que el lagrangiano dependa explícitamente del tiempo.
\end{description}

\subsection{Grados de libertad}

Los vínculos V1--V3 reducen el cuerpo rígido a un solo grado de libertad rotacional. V4--V5
congelan los infinitos grados de libertad del fluido en un único parámetro $h$, y V6 elimina a $h$
como variable dinámica y lo convierte en una función conocida del tiempo. Por tanto
\begin{equation}
\boxed{\;n = 1\;}\qquad\text{grado de libertad.}
\end{equation}

\textbf{Observación importante.} El sistema es \textbf{reónomo}: los vínculos V4--V6 dependen
explícitamente del tiempo. Esto tiene una consecuencia directa que se explota en la
Sección~\ref{sec:conservadas}: la función energía de Jacobi \emph{no} se conserva.

%=======================================================================
\section{Coordenada generalizada y cinemática}
%=======================================================================

\subsection{Elección de la coordenada}

Se adopta como coordenada generalizada el \textbf{ángulo $\theta$ que forma el eje botella--cuerda
con la vertical descendente}:
\begin{equation}
q \equiv \theta, \qquad \theta \in (-\pi,\pi].
\end{equation}
Es la elección natural: es la única variable no fijada por los vínculos, mide directamente la
observable experimental (la inclinación del péndulo en el vídeo) y hace que la posición de
\emph{cualquier} punto del sistema sea función de $\theta$ y de $t$.

\subsection{Geometría}

Medimos todas las distancias desde el pivote $O$ a lo largo del eje de la botella
(Fig.~\ref{fig:esquema}):
\begin{align}
s &= L_c + H && \text{(pivote $\to$ fondo/orificio)},\\
d_b &= L_c + \tfrac{H}{2} && \text{(pivote $\to$ CM de la botella vacía)},\\
d_w(h) &= s - \tfrac{h}{2} && \text{(pivote $\to$ CM de la columna de agua)}.
\end{align}
Nótese que $d_w$ \textbf{aumenta} al vaciarse ($h\downarrow \Rightarrow d_w \uparrow s$): el agua
remanente se concentra cada vez más lejos del pivote. Éste es el mecanismo central del problema.

Con el versor radial $\hat{\mathbf e}_r$ dirigido de $O$ hacia el fondo y $\hat{\mathbf e}_\theta$
perpendicular, la posición de un punto a distancia $r$ del pivote es
\begin{equation}
\mathbf r(r,t) = r\big(\sin\theta,\,-\cos\theta\big), \qquad
\dot{\mathbf r} = \dot r\,\hat{\mathbf e}_r + r\dot\theta\,\hat{\mathbf e}_\theta .
\end{equation}
Para los puntos materiales de la botella $\dot r = 0$; para el centro de masa del conjunto
$\dot r = \dot\ell(t) \neq 0$, porque el CM se desplaza a lo largo del eje conforme se vacía.

\begin{figure}[H]
\centering
\begin{tikzpicture}[scale=0.95,>=Stealth,font=\small]
  \def\ang{20}
  % ---------- techo ----------
  \fill[pattern=north east lines,pattern color=gris] (-2.6,0) rectangle (2.6,0.3);
  \draw[very thick] (-2.6,0)--(2.6,0);
  \coordinate (O) at (0,0);
  \draw[dashed,gris,thin] (O)--(0,-3.2);
  % ---------- cuerpo inclinado ----------
  \begin{scope}[rotate=-\ang]
    \coordinate (T)  at (0,-3.0);   \coordinate (B)  at (0,-6.0);
    \coordinate (CM) at (0,-4.05);
    % cuerda
    \draw[very thick] (O)--(T);
    % agua
    \fill[agua] (-0.72,-6.0)--(0.72,-6.0)--(0.72,-4.55)--(-0.72,-4.55)--cycle;
    \draw[very thick,agua!55!black] (-0.72,-4.55)--(0.72,-4.55);
    % paredes de la botella
    \draw[very thick] (-0.75,-3.0)--(-0.75,-6.0)--(-0.16,-6.0);
    \draw[very thick] (0.75,-3.0)--(0.75,-6.0)--(0.16,-6.0);
    \draw[very thick] (-0.75,-3.0)--(0.75,-3.0);
    % chorro
    \foreach \y in {6.12,6.44,6.76}
      {\draw[agua!80!black,line width=1.4pt] (-0.045,-\y)--(0.045,-\y-0.18);}
    % CM
    \fill[rojo] (CM) circle (2.8pt);
    % cotas (lineas dentro del scope, etiquetas fuera)
    \draw[<->,gris] (-1.15,0)--(-1.15,-3.0);     \coordinate (cLc) at (-1.15,-1.5);
    \draw[<->,gris] (-2.05,0)--(-2.05,-6.0);     \coordinate (cs)  at (-2.05,-3.1);
    \draw[<->,rojo!70] (1.95,0)--(1.95,-4.05);   \coordinate (cl)  at (1.95,-2.1);
    \draw[<->,gris] (1.15,-3.0)--(1.15,-6.0);    \coordinate (cH)  at (1.15,-4.5);
    \draw[<->,agua!55!black] (0.36,-6.0)--(0.36,-4.55); \coordinate (ch) at (0.36,-5.3);
    \draw[<->,gris] (-0.75,-2.68)--(0.75,-2.68); \coordinate (cR) at (0,-2.68);
    \coordinate (cCM) at (-0.12,-4.05);
    \draw[rojo!55,dotted] (0,-4.05)--(1.95,-4.05);
    \coordinate (ca)  at (0.20,-6.05);
    \coordinate (cu)  at (0.22,-6.80);
    \coordinate (cWS) at (-0.78,-4.55);
  \end{scope}
  % ---------- etiquetas (sin rotar) ----------
  \node[left]        at (cLc) {$L_c$};
  \node[left]        at (cs)  {$s$};
  \node[right,rojo]  at (cl)  {$\ell(t)$};
  \node[right]       at (cH)  {$H$};
  \node[right]       at (ch)  {$h$};
  \node[above]       at (cR)  {$2R$};
  \node[left,rojo]   at (cCM) {CM};
  \node[right]       at (ca)  {$a$};
  \node[right,agua!55!black] at (cu) {$u$};
  \node[left,agua!55!black,font=\scriptsize] at (cWS) {sup.\ libre};
  \fill (O) circle (2.3pt);  \node[above left=0pt] at (O) {$O$};
  % angulo
  \draw[->,azulosc,thick] (0,-1.45) arc (-90:-90-\ang:1.45);
  \node[azulosc] at (-0.30,-1.62) {$\theta$};
  % gravedad
  \draw[->,very thick,verde] (-4.15,-0.9)--(-4.15,-1.9);
  \node[verde,left] at (-4.15,-1.4) {$g$};
\end{tikzpicture}
\caption{Geometría del sistema. El pivote $O$ es fijo; $\theta$ es la única coordenada
generalizada. La columna de agua de altura $h(t)$ descansa siempre sobre el fondo, de modo que
su centro de masa se encuentra a $s-h/2$ del pivote y \emph{se aleja} de éste conforme
el nivel baja. El chorro sale por el orificio de área $a$ con velocidad relativa $u$
dirigida a lo largo del eje, es decir, radialmente respecto de $O$.}
\label{fig:esquema}
\end{figure}

%=======================================================================
\section{Propiedades inerciales instantáneas}
\label{sec:inercia}
%=======================================================================

Como la configuración del sistema queda determinada por $h(t)$, todas las propiedades inerciales
son funciones conocidas del tiempo. Definimos la masa de agua
\begin{equation}
m_w(t) = \rho A\,h(t), \qquad A = \pi R^2 .
\end{equation}

\subsection{Masa total y centro de masa}

\begin{align}
M(t) &= m_b + \rho A h , \\[2mm]
\ell(t) &= \frac{m_b d_b + m_w d_w}{M}
        = \frac{m_b\!\left(L_c+\frac H2\right) + \rho A h\!\left(s-\frac h2\right)}{m_b+\rho A h}.
\label{eq:ell}
\end{align}

Es cómodo introducir el \textbf{momento estático} (momento de primer orden respecto del pivote),
que es la cantidad que aparece realmente en las ecuaciones:
\begin{equation}
\boxed{\;S(t) \equiv M(t)\,\ell(t) = m_b d_b + \rho A h\Big(s-\tfrac h2\Big)\;}
\label{eq:S}
\end{equation}

\subsection{Momentos de inercia}

Para un cilindro macizo de masa $m$, radio $R$ y altura $L$, el momento de inercia respecto de un
eje transversal que pasa por su centro de masa es $m(3R^2+L^2)/12 = m(R^2/4 + L^2/12)$. Aplicándolo
a la columna de agua (altura $h$) y usando el teorema de Steiner para trasladar al pivote:
\begin{align}
I_{\text{agua}}^{(O)}(h) &= \rho A h\left[\underbrace{\frac{R^2}{4}+\frac{h^2}{12}}_{\text{propio}}
   + \underbrace{\Big(s-\frac h2\Big)^{\!2}}_{\text{Steiner}}\right], \\[2mm]
I_{\text{bot}}^{(O)} &= I_b + m_b d_b^2 , \qquad
I_b = m_b\left(\frac{R^2}{2}+\frac{H^2}{12}\right)\ \text{(cascarón cilíndrico)} .
\end{align}
El momento de inercia total respecto del pivote es entonces
\begin{equation}
\boxed{\;
I_O(t) = I_b + m_b d_b^2 + \rho A h\left[\frac{R^2}{4}+\frac{h^2}{12}+\Big(s-\frac h2\Big)^{\!2}\right]
\;}
\label{eq:IO}
\end{equation}
y el momento de inercia respecto del centro de masa instantáneo, por Steiner,
\begin{equation}
I_{\rm cm}(t) = I_O(t) - M(t)\,\ell(t)^2 = I_O - \frac{S^2}{M}.
\label{eq:Icm}
\end{equation}

\subsection{Dos identidades clave}
\label{sec:identidades}

Derivando \eqref{eq:S} y \eqref{eq:IO} respecto de $h$ se obtienen dos expresiones que serán
decisivas más adelante:
\begin{equation}
\boxed{\;\frac{\dd S}{\dd h}=\rho A\,(s-h),\qquad
        \frac{\dd I_O}{\dd h}=\rho A\left[(s-h)^2+\frac{R^2}{4}\right].\;}
\label{eq:identidades}
\end{equation}

\textbf{Interpretación.} Ambas son exactamente el momento estático y el momento de inercia de una
rodaja de espesor $\dd h$ situada \emph{en la superficie libre}, a distancia $s-h$ del pivote.
Es decir: desde el punto de vista geométrico, al vaciarse el tanque el sistema ``pierde'' masa
\emph{por arriba} (la superficie desciende), aunque el agua salga \emph{por abajo}.
Esta asimetría entre dónde desaparece la masa en la descripción y dónde sale físicamente es el
origen del término disipativo que se obtendrá en la Sección~\ref{sec:masavar}.

%=======================================================================
\section{Ley de cambio de masa: vaciado de un tanque cilíndrico}
\label{sec:torricelli}
%=======================================================================

\subsection{Deducción de la ley de Torricelli}

Apliquemos Bernoulli a lo largo de una línea de corriente entre la superficie libre (punto 1,
presión atmosférica $p_0$, velocidad $v_1$, altura $h$ sobre el orificio) y la vena contracta del
chorro (punto 2, presión $p_0$, velocidad $u$, altura $0$):
\begin{equation}
p_0 + \tfrac12\rho v_1^2 + \rho g h = p_0 + \tfrac12 \rho u^2 .
\end{equation}
La continuidad impone $A v_1 = a\,u$, de donde
\begin{equation}
u = \sqrt{\frac{2gh}{1-(a/A)^2}} \;\simeq\; \sqrt{2gh}
\qquad\text{para } a \ll A .
\end{equation}
Para el caso numérico de la Sección~\ref{sec:numerico} se tiene $a/A = 3.9\times10^{-3}$, luego la
corrección es de $8\times10^{-6}$ y es perfectamente lícito quedarse con la forma clásica.
Introduciendo el coeficiente de descarga $C_d$ (contracción de la vena y pérdidas viscosas;
$C_d\approx 0.60$--$0.62$ para un orificio circular de pared delgada), el gasto volumétrico es
\begin{equation}
Q = C_d\,a\,\sqrt{2gh}.
\end{equation}

\subsection{Integración: la ley $h(t)$}

La conservación de la masa en el tanque cilíndrico, $A\,\dd h = -Q\,\dd t$, da la ecuación
diferencial separable
\begin{equation}
\frac{\dd h}{\dd t} = -\frac{C_d\,a}{A}\sqrt{2g\,h}\,.
\label{eq:edoh}
\end{equation}
Separando variables e integrando desde $h(0)=h_0$:
\begin{equation}
\int_{h_0}^{h}\frac{\dd h'}{\sqrt{h'}} = -\frac{C_d a}{A}\sqrt{2g}\int_0^t \dd t'
\;\Longrightarrow\;
2\left(\sqrt h - \sqrt{h_0}\right) = -\frac{C_d a}{A}\sqrt{2g}\;t .
\end{equation}
Despejando y definiendo el \textbf{tiempo de vaciado}
\begin{equation}
\boxed{\;t_v = \frac{A}{C_d\,a}\sqrt{\frac{2h_0}{g}}\;}
\label{eq:tv}
\end{equation}
se obtiene la ley de vaciado, que es \emph{cuadrática} en el tiempo:
\begin{equation}
\boxed{\;h(t) = h_0\left(1-\frac{t}{t_v}\right)^{\!2},\qquad 0\le t\le t_v.\;}
\label{eq:ht}
\end{equation}

\subsection{Función de cambio de masa}

De \eqref{eq:ht} se sigue inmediatamente
\begin{align}
m_w(t) &= \rho A h_0\left(1-\frac{t}{t_v}\right)^{\!2}
       = m_{w0}\left(1-\frac{t}{t_v}\right)^{\!2},
\label{eq:mt}\\[2mm]
\dot m_w(t) &= -\frac{2\,m_{w0}}{t_v}\left(1-\frac{t}{t_v}\right)
            = -\rho\,C_d\,a\sqrt{2g\,h(t)}\;\le 0 .
\label{eq:mdot}
\end{align}
Denotaremos por $\dot m_{\text{out}} \equiv -\dot m_w = \rho C_d a\sqrt{2gh} \ge 0$ el gasto
másico que abandona el sistema. Obsérvese que $\dot m_{\text{out}}$ es \emph{máximo al principio}
y tiende linealmente a cero al final: el vaciado es fuertemente no uniforme, lo que descarta el
modelo de gasto constante.

\subsection{Corrección por gravedad efectiva (semilla del efecto paramétrico)}
\label{sec:geff}

En rigor, la presión que impulsa el chorro no es $\rho g h$ sino $\rho\,g_{\rm eff}h$, donde
$g_{\rm eff}$ es la componente axial del campo de gravedad efectivo en el marco no inercial de la
botella. Para el péndulo,
\begin{equation}
g_{\rm eff} \simeq g\cos\theta + s\,\dot\theta^2 ,
\label{eq:geff}
\end{equation}
donde el segundo término es la contribución centrífuga (dirigida del pivote hacia el fondo, por
tanto se \emph{suma}). Con $\theta=\Theta\cos\omega t$ y $\omega^2 = g/L_{\rm eq}$,
\begin{equation}
\frac{g_{\rm eff}}{g} \simeq 1 + \Theta^2\!\left(\frac{s}{2L_{\rm eq}}-\frac14\right)
 - \Theta^2\!\left(\frac{s}{2L_{\rm eq}}+\frac14\right)\cos 2\omega t .
\label{eq:geffmod}
\end{equation}
Aparece así una \textbf{modulación del gasto a la frecuencia $2\omega$}, de orden $\Theta^2$, que
es exactamente la estructura necesaria para una resonancia paramétrica. Su magnitud se evalúa en
la Sección~\ref{sec:parametrico}; se anticipa que es demasiado pequeña para producir crecimiento
en el régimen lento, y se determina bajo qué condiciones dejaría de serlo.

%=======================================================================
\section{El lagrangiano del sistema}
\label{sec:lagrangiano}
%=======================================================================

\subsection{Construcción término a término}

Siguiendo el requisito del enunciado, construimos $T$ como suma de la energía cinética de
traslación del centro de masa y la energía cinética de rotación en torno a él (teorema de König),
y $U$ como la energía potencial gravitatoria del centro de masa.

\paragraph{(i) Energía cinética del centro de masa.}
El centro de masa se halla a distancia $\ell(t)$ del pivote sobre el eje, de modo que
$\mathbf r_{\rm cm} = \ell(t)\,\hat{\mathbf e}_r(\theta)$ y
\begin{equation}
\mathbf v_{\rm cm} = \dot\ell\,\hat{\mathbf e}_r + \ell\dot\theta\,\hat{\mathbf e}_\theta
\;\Longrightarrow\;
v_{\rm cm}^2 = \dot\ell^{\,2} + \ell^2\dot\theta^2 .
\end{equation}
Por tanto
\begin{equation}
T_{\rm cm} = \tfrac12 M(t)\left[\dot\ell(t)^2 + \ell(t)^2\dot\theta^2\right].
\end{equation}
El término $\tfrac12 M\dot\ell^{\,2}$ representa la energía cinética asociada al
\emph{desplazamiento longitudinal} del centro de masa durante el vaciado. Es una función
\textbf{exclusivamente del tiempo}, $f(t)=\tfrac12 M(t)\dot\ell(t)^2$, y por tanto satisface
$\partial f/\partial\theta = \partial f/\partial\dot\theta = 0$: no contribuye a la ecuación de
Euler--Lagrange. Puede eliminarse mediante la transformación de calibre
$\mathcal L \to \mathcal L - \dd F/\dd t$ con $F(t)=\int^t f$. Lo conservamos escrito para dejar
constancia de que se ha tenido en cuenta.

\paragraph{(ii) Energía cinética de rotación.}
El sólido (botella $+$ agua congelada) gira con velocidad angular $\dot\theta\,\hat{\mathbf k}$
alrededor del eje perpendicular al plano que pasa por el centro de masa:
\begin{equation}
T_{\rm rot} = \tfrac12 I_{\rm cm}(t)\,\dot\theta^2,
\qquad I_{\rm cm} = I_O - \frac{S^2}{M}.
\end{equation}

\paragraph{(iii) Energía potencial gravitatoria.}
Tomando el pivote como origen de potenciales, la altura del centro de masa es
$y_{\rm cm} = -\ell(t)\cos\theta$, de donde
\begin{equation}
U = -M(t)\,g\,\ell(t)\cos\theta = -g\,S(t)\cos\theta .
\end{equation}

\subsection{El lagrangiano}

Reuniendo los tres términos:
\begin{equation}
\boxed{\;
\mathcal L(\theta,\dot\theta,t) =
\underbrace{\tfrac12 M(t)\!\left[\dot\ell^{\,2}+\ell^2\dot\theta^2\right]}_{T_{\rm cm}}
+\underbrace{\tfrac12 I_{\rm cm}(t)\dot\theta^2}_{T_{\rm rot}}
+\underbrace{g\,S(t)\cos\theta}_{-U}\;}
\label{eq:Lfull}
\end{equation}

Usando \eqref{eq:Icm}, los dos términos cuadráticos en $\dot\theta$ se combinan:
$M\ell^2 + I_{\rm cm} = M\ell^2 + I_O - M\ell^2 = I_O$. Descartando el término de calibre,
el lagrangiano adopta la forma compacta de un \textbf{péndulo físico de parámetros variables}:
\begin{equation}
\boxed{\;
\mathcal L(\theta,\dot\theta,t) = \tfrac12\,I_O(t)\,\dot\theta^2 + g\,S(t)\,\cos\theta\;}
\label{eq:L}
\end{equation}
con $I_O(t)$ y $S(t)$ dados por \eqref{eq:IO} y \eqref{eq:S} evaluados en $h(t)$ de \eqref{eq:ht}.
Obsérvese que \emph{toda} la física del vaciado ha quedado encapsulada en dos funciones escalares
del tiempo.

\subsection{Ecuación de Euler--Lagrange}

Con $q=\theta$:
\begin{equation}
\frac{\partial \mathcal L}{\partial \dot\theta} = I_O(t)\,\dot\theta \equiv p_\theta ,
\qquad
\frac{\partial \mathcal L}{\partial \theta} = -g\,S(t)\sin\theta .
\end{equation}
La ecuación $\dfrac{\dd}{\dd t}\!\left(\dfrac{\partial\mathcal L}{\partial\dot\theta}\right)
-\dfrac{\partial\mathcal L}{\partial\theta}=0$ da
\begin{equation}
\frac{\dd}{\dd t}\Big(I_O\dot\theta\Big) + g\,S\sin\theta = 0
\qquad\Longleftrightarrow\qquad
I_O\ddot\theta + \dot I_O\,\dot\theta + g\,S\sin\theta = 0 .
\label{eq:ELnaive}
\end{equation}

El término $\dot I_O\dot\theta$ aparece \emph{automáticamente}: es la firma de la variabilidad
inercial. Sin embargo, como se demuestra en la Sección~\ref{sec:masavar}, la ecuación
\eqref{eq:ELnaive} es \textbf{incorrecta} para este sistema, porque el formalismo lagrangiano
estándar presupone un conjunto \emph{fijo} de partículas. La corrección no es un detalle:
cambia la predicción de amplitud final en un factor $\sim 8$.

%=======================================================================
\section{Coordenadas cíclicas y cantidades conservadas}
\label{sec:conservadas}
%=======================================================================

\subsection{¿Hay coordenadas cíclicas?}

Una coordenada $q_k$ es cíclica si $\partial\mathcal L/\partial q_k=0$. En \eqref{eq:L},
\begin{equation}
\frac{\partial\mathcal L}{\partial\theta} = -gS(t)\sin\theta \neq 0 ,
\end{equation}
luego \textbf{$\theta$ no es cíclica} y el momento conjugado $p_\theta = I_O\dot\theta$
(el momento angular respecto del pivote) \textbf{no se conserva}. El responsable es el torque
gravitatorio: es el único grado de libertad y está acoplado a la gravedad.

Existen, no obstante, dos situaciones límite instructivas:

\begin{itemize}[leftmargin=*]
\item \textbf{Límite $g\to 0$ (o pivote en caída libre).} Entonces $\mathcal L=\tfrac12 I_O\dot\theta^2$,
$\theta$ es cíclica y $p_\theta=I_O\dot\theta=$ cte. El sistema sería un análogo rotacional exacto
de la ``patinadora'' que abre y cierra los brazos: al vaciarse, $I_O$ disminuye y $\dot\theta$
aumenta. (Con la corrección de masa variable de la Sec.~\ref{sec:masavar} lo conservado pasa a ser
$p_\theta$ menos el momento angular evacuado por el chorro.)

\item \textbf{Generalización tridimensional.} Si se relaja el vínculo V3 y se permite el movimiento
azimutal (péndulo esférico), con coordenadas $(\theta,\varphi)$,
\begin{equation}
\mathcal L = \tfrac12 I_O(t)\left(\dot\theta^2+\sin^2\!\theta\,\dot\varphi^2\right) + gS(t)\cos\theta ,
\end{equation}
la coordenada $\varphi$ \emph{sí} es cíclica ($\partial\mathcal L/\partial\varphi = 0$, por la
simetría de revolución del campo gravitatorio y de la botella cilíndrica en torno a la vertical).
El momento conjugado
\begin{equation}
p_\varphi = I_O(t)\,\sin^2\!\theta\;\dot\varphi
\end{equation}
sería constante en el modelo ingenuo. Sin embargo, el chorro se lleva momento angular vertical, de
modo que la ley correcta es
$\dot p_\varphi = -\dot m_{\text{out}}\,s^2\sin^2\!\theta\,\dot\varphi$: la simetría existe, pero
el sistema es abierto y el teorema de Noether no basta. Ésta es una advertencia general para
sistemas de masa variable: \emph{una coordenada cíclica no garantiza una cantidad conservada si el
sistema intercambia materia con el exterior}.
\end{itemize}

\subsection{La función energía de Jacobi}

Para un lagrangiano $\mathcal L(q,\dot q,t)$ la función energía (integral de Jacobi) es
\begin{equation}
h_J \equiv \dot\theta\,\frac{\partial\mathcal L}{\partial\dot\theta}-\mathcal L
= \tfrac12 I_O(t)\dot\theta^2 - g S(t)\cos\theta .
\end{equation}
Como $T$ es una forma cuadrática homogénea en $\dot\theta$, $h_J$ coincide numéricamente con la
energía mecánica $E=T+U$. Su evolución obedece el teorema general
\begin{equation}
\frac{\dd h_J}{\dd t} = -\frac{\partial\mathcal L}{\partial t},
\end{equation}
que aquí da explícitamente
\begin{equation}
\boxed{\;
\frac{\dd E}{\dd t} = -\frac{\partial\mathcal L}{\partial t}
= -\tfrac12 \dot I_O(t)\,\dot\theta^2 - g\,\dot S(t)\cos\theta \;\neq 0 .}
\label{eq:dEdt}
\end{equation}

\textbf{Conclusión:} \emph{no existe ninguna cantidad exactamente conservada}. El sistema es
reónomo (los vínculos V4--V6 dependen de $t$) y abierto (pierde materia). Usando
\eqref{eq:identidades} y $\dot h<0$, ambos términos de \eqref{eq:dEdt} son positivos para
$|\theta|<\pi/2$: la energía \emph{aumenta} formalmente, pero esto no es paradójico: sólo refleja
que estamos siguiendo un sistema cuya definición material cambia. El agua que sale se lleva su
propia energía, que ya no se contabiliza.

\subsection{Lo que sí se conserva: el invariante adiabático}
\label{sec:adiabatico}

Cuando el vaciado es lento comparado con el período, existe una cantidad
\emph{aproximadamente} conservada mucho más útil que $E$: la \textbf{acción}
\begin{equation}
\boxed{\;
J = \oint p_\theta\,\dd\theta = \pi\, I_O(t)\,\omega(t)\,\Theta(t)^2
  = \pi\,\Theta^2\sqrt{g\,S(t)\,I_O(t)} \;}
\label{eq:accion}
\end{equation}
donde $\Theta$ es la amplitud y $\omega=\sqrt{gS/I_O}$ la frecuencia instantánea.
El teorema del invariante adiabático garantiza que $J$ es constante hasta correcciones
exponencialmente pequeñas en el parámetro
\begin{equation}
\eps \equiv \frac{|\dot\omega|}{\omega^2} \ll 1 .
\label{eq:epsilon}
\end{equation}
Para el oscilador \emph{hamiltoniano} de inercia variable
$\frac{\dd}{\dd t}(I\dot\theta)+K\theta=0$, la conservación de $J$ es exacta al orden adiabático.
Como se demuestra en la Sección~\ref{sec:masavar}, la pérdida de masa por el orificio rompe esta
conservación de una forma perfectamente calculable:
\begin{equation}
\boxed{\;\frac{\dd J}{\dd t} = -\frac{\dot m_{\text{out}}\,s^2}{I_O}\;J \;}
\label{eq:dJdt}
\end{equation}
Es decir: \emph{el invariante adiabático decae exactamente al ritmo con que el chorro evacua
momento angular}. Esta ecuación es el resultado teórico central del trabajo y es la que se
confronta con el experimento.

%=======================================================================
\section{Sistemas de masa variable: corrección del lagrangiano}
\label{sec:masavar}
%=======================================================================

\subsection{Por qué \eqref{eq:ELnaive} es incompleta}

El principio de Hamilton se formula para un conjunto \emph{fijo} de partículas. Cuando se escribe
$\mathcal L = \tfrac12 I_O(t)\dot\theta^2 + gS(t)\cos\theta$ y se sustituye $I_O(t)$ ``a mano'',
se está describiendo a cada instante un sistema material \emph{distinto}, sin decir nada sobre
qué le ocurre a la masa que desaparece. La ecuación de Euler--Lagrange resultante supone
implícitamente que la masa se evapora \emph{en reposo relativo respecto del punto donde estaba}, y
en la posición donde la geometría dice que desaparece, es decir, \textbf{en la superficie libre}.
Pero el agua no sale por la superficie libre: sale por el \textbf{fondo}, a distancia $s$ del
pivote, no a distancia $s-h$. Ésa es la discrepancia, y como veremos es cuantitativamente enorme.

\subsection{Balance de momento angular con flujo de masa}

Tomamos como volumen de control $\mathcal V$ el interior de la botella (que se mueve y se deforma
con ella) y aplicamos el teorema de transporte de Reynolds al momento angular respecto del pivote
fijo $O$:
\begin{equation}
\frac{\dd}{\dd t}\int_{\mathcal V}\rho\,(\mathbf r\times\mathbf v)\,\dd V
\;+\;\oint_{\partial\mathcal V}\rho\,(\mathbf r\times\mathbf v)\,
(\mathbf v-\mathbf v_{\partial\mathcal V})\cdot\hat{\mathbf n}\,\dd \Sigma
\;=\;\boldsymbol\tau_{\rm ext}.
\label{eq:reynolds}
\end{equation}

Analicemos cada término.

\paragraph{Superficies sin flujo.} En las paredes y en la superficie libre se cumple
$(\mathbf v-\mathbf v_{\partial\mathcal V})\cdot\hat{\mathbf n}=0$: las paredes son impermeables y
la superficie libre es una superficie material que desciende con el fluido. No contribuyen.

\paragraph{Flujo por el orificio.} La velocidad del fluido a la salida es la del orificio más la
velocidad relativa: $\mathbf v = s\dot\theta\,\hat{\mathbf e}_\theta - u\,\hat{\mathbf e}_r$.
Como el orificio está \emph{sobre el eje} que pasa por el pivote, $\mathbf r = s\,\hat{\mathbf e}_r$
y por tanto
\begin{equation}
\mathbf r\times\mathbf v = s\hat{\mathbf e}_r\times
\big(s\dot\theta\,\hat{\mathbf e}_\theta - u\,\hat{\mathbf e}_r\big) = s^2\dot\theta\,\hat{\mathbf k}.
\end{equation}

\begin{resultado}[Velocidad relativa del chorro]
La velocidad relativa de salida $u$ \textbf{no aporta momento angular} respecto del pivote, porque
es paralela a $\mathbf r$. El chorro no ejerce torque sobre el sistema. Pero sí \emph{extrae}
momento angular, a razón de $\dot m_{\text{out}}\,s^2\dot\theta$ por unidad de tiempo:
la masa que se va se lleva su parte de momento angular de arrastre. Éste es el sentido preciso en
que ``el agua sale con velocidad relativa'': la elección importa para el balance de momento
angular, no a través de un empuje, sino a través del \emph{brazo} $s$ al que se produce la
evacuación.
\end{resultado}

\paragraph{Momento angular contenido.} Bajo la hipótesis de agua congelada (I5), la contribución
de la deriva axial a $\mathbf r\times\mathbf v$ se anula por simetría de revolución, de modo que
\begin{equation}
\int_{\mathcal V}\rho\,(\mathbf r\times\mathbf v)\,\dd V = I^{(O)}_{\text{agua}}(h)\;\dot\theta\,\hat{\mathbf k},
\end{equation}
y análogamente para la botella. Sumando ambas y usando \eqref{eq:identidades}:
\begin{equation}
\frac{\dd}{\dd t}\Big(I_O\dot\theta\Big) + \dot m_{\text{out}}\,s^2\,\dot\theta = -g\,S\,\sin\theta .
\label{eq:balance}
\end{equation}

\subsection{La ecuación de movimiento correcta}

Desarrollando la derivada y agrupando obtenemos el resultado central:
\begin{equation}
\boxed{\;
I_O(t)\,\ddot\theta \;+\; \Gamma(t)\,\dot\theta \;+\; g\,S(t)\,\sin\theta \;=\; 0
\;}
\label{eq:EOM}
\end{equation}
con el \textbf{coeficiente de frenado por vaciado}
\begin{equation}
\Gamma(t) \;\equiv\; \dot I_O + \dot m_{\text{out}}\,s^2 .
\label{eq:Gammadef}
\end{equation}
Usando \eqref{eq:identidades} y $\dot m_{\text{out}}=-\rho A\dot h$:
\begin{equation}
\Gamma = \rho A\dot h\left[(s-h)^2+\frac{R^2}{4}\right] - \rho A\dot h\,s^2
       = \rho A\,\dot h\left[\frac{R^2}{4}-h\,(2s-h)\right],
\end{equation}
y como $\dot h<0$, escribiendo $|\dot h|=-\dot h$:
\begin{equation}
\boxed{\;\Gamma(t) = \rho A\,|\dot h|\left[h\,(2s-h)-\frac{R^2}{4}\right] \;>\;0 .\;}
\label{eq:Gamma}
\end{equation}

\begin{resultado}[Signo del frenado]
$\Gamma>0$ siempre que $h(2s-h)>R^2/4$, es decir, para
$h > h_c \simeq R^2/(8s)$. Con los valores de la Sec.~\ref{sec:numerico},
$h_c=\SI{0.27}{\milli\meter}$, correspondiente a \SI{1.3}{\gram} de agua residual.
Salvo en esa última película de líquido, \textbf{el vaciado siempre frena la oscilación}.
\end{resultado}

\paragraph{Interpretación física.} El factor entre corchetes es $s^2-(s-h)^2-R^2/4$: la diferencia
entre el momento de inercia específico del punto por donde el agua \emph{sale} ($s^2$) y el del
punto por donde la geometría dice que \emph{desaparece} ($\,(s-h)^2+R^2/4$). El sistema evacua más
momento angular del que ``libera'' al encoger, y el déficit lo paga el agua que queda: de ahí el
frenado. Cuanto más alta es la columna (mayor $h$), mayor es la separación entre la superficie
libre y el orificio, y mayor es el frenado.

\subsection{Control: el límite de masa puntual}
\label{sec:puntual}

Si toda la masa estuviera concentrada en el orificio ($h\to 0$ con $m_w$ fija, o una botella
puntual a distancia $s$), entonces $I_O=Ms^2$, $S=Ms$ y
\begin{equation}
\dot I_O = \dot M s^2 = -\dot m_{\text{out}}s^2
\quad\Longrightarrow\quad \Gamma \equiv 0 ,
\end{equation}
y la ecuación \eqref{eq:EOM} se reduce a
\begin{equation}
M s^2\ddot\theta + Mgs\sin\theta = 0
\quad\Longleftrightarrow\quad
\ddot\theta + \frac{g}{s}\sin\theta = 0 .
\end{equation}

\begin{resultado}[El modelo de masa puntual predice \emph{nada}]
Para un péndulo simple que pierde masa sin cambiar de geometría, \textbf{el período y la amplitud
permanecen exactamente constantes}, sea cual sea la ley $m(t)$. La masa se cancela idénticamente.
Todo el fenómeno observado en el laboratorio es, por tanto, un efecto de \emph{cuerpo extenso}:
depende de que la distribución de masa cambie de \emph{forma}, no de que cambie de \emph{valor}.
Esto justifica de manera contundente la elección del modelo de péndulo físico.
\end{resultado}

\subsection{Comparación con el lagrangiano ingenuo}

El lagrangiano ingenuo \eqref{eq:ELnaive} equivale a hacer $\Gamma\to\dot I_O$ en \eqref{eq:EOM}.
La diferencia entre ambos modelos es precisamente el término de flujo:
\begin{equation}
\Gamma - \dot I_O = \dot m_{\text{out}}\,s^2 .
\end{equation}
Como $\dot I_O<0$ (el sistema pierde inercia), el modelo ingenuo contiene un término
\emph{antiamortiguador} que bombea energía a la oscilación. Veremos en la
Sección~\ref{sec:numerico} que predice una amplitud final de \ang{71} partiendo de \ang{12}:
un resultado manifiestamente falso que cualquier observación descarta de inmediato. La
comparación entre ambos modelos constituye, de hecho, un excelente test experimental.

%=======================================================================
\section{Solución de la ecuación de Euler--Lagrange}
\label{sec:solucion}
%=======================================================================

\subsection{Frecuencia instantánea y longitud equivalente}

Linealizando \eqref{eq:EOM} para $|\theta|\ll 1$:
\begin{equation}
\ddot\theta + \gamma(t)\,\dot\theta + \omega^2(t)\,\theta = 0,
\qquad
\gamma \equiv \frac{\Gamma}{I_O},\qquad
\omega^2(t) \equiv \frac{g\,S(t)}{I_O(t)} .
\label{eq:lineal}
\end{equation}
Es cómodo introducir la \textbf{longitud equivalente} (longitud del péndulo simple isócrono):
\begin{equation}
\boxed{\;
L_{\rm eq}(t) = \frac{I_O(t)}{S(t)} = \frac{I_O}{M\ell},
\qquad
\omega^2 = \frac{g}{L_{\rm eq}},\qquad
T(t) = 2\pi\sqrt{\frac{L_{\rm eq}(t)}{g}} \;}
\label{eq:Leq}
\end{equation}
Explícitamente, sustituyendo \eqref{eq:S} y \eqref{eq:IO} con $h=h_0(1-t/t_v)^2$:
\begin{equation}
L_{\rm eq}(t) =
\frac{\,I_b + m_b d_b^2 + \rho A h\left[\frac{R^2}{4}+\frac{h^2}{12}+\left(s-\frac h2\right)^2\right]}
     {\,m_b d_b + \rho A h\left(s-\frac h2\right)} .
\label{eq:Leqexp}
\end{equation}

$L_{\rm eq}$ tiene un significado geométrico preciso: el punto situado a distancia $L_{\rm eq}$ del
pivote es el \textbf{centro de oscilación} (o centro de percusión) del péndulo compuesto.

\subsection{El período es no monótono: condición del extremo}
\label{sec:extremo}

Buscamos los extremos de $L_{\rm eq}=I_O/S$ como función de $h$. Derivando,
\begin{equation}
\frac{\dd L_{\rm eq}}{\dd h} = \frac{I_O' S - I_O S'}{S^2} = 0
\quad\Longleftrightarrow\quad
\frac{I_O'}{S'} = \frac{I_O}{S}=L_{\rm eq},
\end{equation}
y sustituyendo las identidades \eqref{eq:identidades}:
\begin{equation}
\boxed{\;
\frac{(s-h^\ast)^2 + R^2/4}{\,s-h^\ast\,} = L_{\rm eq}(h^\ast)
\qquad\Longrightarrow\qquad
s-h^\ast \simeq L_{\rm eq}(h^\ast)\;}
\label{eq:extremo}
\end{equation}

\begin{resultado}[Criterio del centro de oscilación]
\emph{El período alcanza su valor extremo en el instante en que la superficie libre del agua pasa
exactamente por el centro de oscilación instantáneo del péndulo.}
\end{resultado}

La interpretación es transparente y es el resultado cualitativo más bonito del problema. Es un
hecho conocido de la teoría del péndulo compuesto que \emph{añadir} masa a una distancia $r$ del
pivote aumenta $L_{\rm eq}$ si $r>L_{\rm eq}$ y lo disminuye si $r<L_{\rm eq}$. Aquí ocurre lo
contrario porque se \emph{quita} masa, y siempre desde la superficie libre, situada a $r=s-h$:

\begin{itemize}[leftmargin=*]
\item \textbf{Al principio} la botella está llena, la superficie está \emph{arriba}
($r=s-H$, cerca del pivote), luego $r<L_{\rm eq}$: quitar masa de ahí \textbf{aumenta}
$L_{\rm eq}$ y el período \textbf{crece}.
\item \textbf{Conforme baja el nivel}, $r=s-h$ aumenta. Cuando cruza $L_{\rm eq}$ se alcanza el
máximo del período.
\item \textbf{Al final}, la superficie está \emph{debajo} del centro de oscilación
($r>L_{\rm eq}$): quitar masa de ahí \textbf{disminuye} $L_{\rm eq}$ y el período \textbf{decrece}.
\item \textbf{Al vaciarse por completo} el sistema vuelve a ser la botella vacía, luego
$T(t_v)=T_{\rm vacía}$, un valor fijado sólo por la geometría del recipiente. Si la botella vacía
y la llena tienen distribuciones de masa geométricamente semejantes (ambas llenan el mismo cilindro
de manera uniforme), entonces $T_{\rm vacía}\simeq T_{\rm llena}$ y \textbf{la curva $T(t)$ empieza
y termina en el mismo valor}, con una joroba en medio.
\end{itemize}

\subsection{Solución WKB: la ley de la amplitud}
\label{sec:wkb}

La ecuación \eqref{eq:lineal} tiene coeficientes lentamente variables. Eliminamos primero el
término de primer orden mediante la sustitución estándar
\begin{equation}
\theta(t) = \psi(t)\,\exp\left[-\frac12\int_0^t\gamma(t')\,\dd t'\right],
\end{equation}
que transforma \eqref{eq:lineal} en
\begin{equation}
\ddot\psi + \left(\omega^2 - \frac{\gamma^2}{4} - \frac{\dot\gamma}{2}\right)\psi = 0
\;\simeq\; \ddot\psi + \omega^2(t)\,\psi = 0 ,
\end{equation}
donde se han despreciado $\gamma^2$ y $\dot\gamma$ frente a $\omega^2$ (para los valores de la
Sec.~\ref{sec:numerico}, $\gamma/\omega \lesssim 4\times10^{-3}$).
La solución WKB de esta ecuación es
$\psi \propto \omega^{-1/2}\cos\!\big(\int_0^t\omega\,\dd t'+\phi_0\big)$, luego
\begin{equation}
\boxed{\;
\theta(t) = \Theta(t)\,\cos\!\left(\int_0^t\!\omega(t')\dd t' + \phi_0\right),
\qquad
\Theta(t) = \Theta_0\sqrt{\frac{\omega_0}{\omega(t)}}\;
\exp\!\left[-\frac12\int_0^t\!\frac{\Gamma}{I_O}\dd t'\right].}
\label{eq:wkb}
\end{equation}

\subsection{Descomposición en dos mecanismos}

La expresión \eqref{eq:wkb} se ilumina al separar $\Gamma=\dot I_O + \dot m_{\text{out}}s^2$.
Usando $\exp\left[-\frac12\int \dot I_O/I_O\right]=\sqrt{I_0/I_O}$ y $\omega=\sqrt{gS/I_O}$:
\begin{equation}
\boxed{\;
\frac{\Theta(t)}{\Theta_0} =
\underbrace{\left[\frac{S_0\,I_{O,0}}{S(t)\,I_O(t)}\right]^{1/4}}_{\text{(A) bombeo adiabático}}
\times
\underbrace{\exp\left[-\frac{s^2}{2}\int_0^t \frac{\dot m_{\text{out}}}{I_O}\,\dd t'\right]}
_{\text{(B) frenado por el chorro}} \;}
\label{eq:amplitud}
\end{equation}

\paragraph{(A) Bombeo adiabático.} Es exactamente la conservación de la acción
\eqref{eq:accion}: $J=\pi\Theta^2\sqrt{gSI_O}=$ cte $\Rightarrow \Theta\propto (S I_O)^{-1/4}$.
Como el vaciado reduce tanto $S$ como $I_O$, este factor \textbf{siempre aumenta} la amplitud. Es
el mismo mecanismo por el que un péndulo cuya cuerda se acorta lentamente oscila con mayor ángulo.

\paragraph{(B) Frenado por el chorro.} Es la violación del invariante adiabático anunciada en
\eqref{eq:dJdt}: el agua se lleva momento angular al brazo máximo $s$. \textbf{Siempre reduce} la
amplitud.

\begin{resultado}[Competencia]
La evolución de la amplitud es el resultado de una competencia entre dos efectos que pueden ser
\emph{individualmente muy grandes} y de signo opuesto. En el caso numérico de la
Sec.~\ref{sec:numerico}, el factor (A) vale $\times 5.69$ y el factor (B) vale $\times 0.127$; su
producto es $\times 0.72$. Un modelo que omita cualquiera de los dos se equivoca por casi un orden
de magnitud.
\end{resultado}

\subsection{Universalidad: el resultado experimental más limpio}
\label{sec:universalidad}

Definamos la integral de frenado acumulada
\begin{equation}
\mathcal K \equiv \int_0^{t_v}\frac{\Gamma}{I_O}\,\dd t .
\end{equation}
Cambiando la variable de integración de $t$ a $h$ mediante $\dd t = \dd h/\dot h$, el factor
$\dot h$ contenido en $\Gamma$ \emph{se cancela idénticamente}:
\begin{equation}
\boxed{\;
\mathcal K = \rho A \int_0^{h_0}
\frac{h\,(2s-h)-R^2/4}{\;I_O(h)\;}\;\dd h \;}
\label{eq:K}
\end{equation}

\begin{resultado}[Independencia del tamaño del orificio]
$\mathcal K$ es una integral \textbf{puramente geométrica}: no contiene $a$, ni $C_d$, ni $g$, ni
$t_v$. En consecuencia, el factor total de cambio de amplitud a lo largo de un vaciado completo,
\begin{equation}
\frac{\Theta(t_v)}{\Theta_0} = \sqrt{\frac{\omega_0}{\omega_f}}\;e^{-\mathcal K/2},
\end{equation}
\emph{no depende de la rapidez con que se vacíe la botella}. Perforar un agujero el doble de grande
hace que todo ocurra cuatro veces más rápido, pero la amplitud final es la misma.
Ésta es una predicción falsable, robusta y fácil de contrastar: basta repetir el experimento con
orificios de distintos diámetros.
\end{resultado}

\subsection{Corrección de amplitud finita}

Las oscilaciones reales no son infinitesimales. La corrección clásica al período de un péndulo
(desarrollo de la integral elíptica completa) se aplica aquí con $L_{\rm eq}$ instantánea:
\begin{equation}
\boxed{\;
T(t) = 2\pi\sqrt{\frac{L_{\rm eq}(t)}{g}}
\left[1+\frac{\Theta(t)^2}{16}+\frac{11\,\Theta(t)^4}{3072}+\dots\right]}
\label{eq:Tfinita}
\end{equation}
Este término acopla las dos preguntas del enunciado: como la amplitud cambia, el período adquiere
una contribución adicional dependiente del tiempo. Para $\Theta=\ang{12}$ la corrección vale
$+0.27\%$, comparable a la variación del período que buscamos medir ($+6\%$) y desde luego mayor
que la precisión típica de un cronometraje por vídeo. \textbf{Debe incluirse al comparar con el
experimento.}

\subsection{Resumen: el sistema de ecuaciones a resolver}

\begin{equation}
\left\{
\begin{aligned}
h(t) &= h_0(1-t/t_v)^2, \qquad t_v=\frac{A}{C_d a}\sqrt{2h_0/g},\\[1mm]
S(t) &= m_b d_b + \rho A h\,(s-h/2),\\[1mm]
I_O(t) &= I_b+m_b d_b^2+\rho A h\big[R^2/4+h^2/12+(s-h/2)^2\big],\\[1mm]
\Gamma(t) &= \rho A|\dot h|\big[h(2s-h)-R^2/4\big],\\[1mm]
0 &= I_O\ddot\theta + \Gamma\dot\theta + gS\sin\theta .
\end{aligned}\right.
\label{eq:sistema}
\end{equation}

Todo está determinado por la geometría $(H,R,L_c,a,C_d)$ y las masas $(m_b,\rho)$. No hay
parámetros libres de ajuste salvo $C_d$, que además sólo fija la escala temporal $t_v$ y no afecta
a las amplitudes.

%=======================================================================
\section{Resultados numéricos}
\label{sec:numerico}
%=======================================================================

\subsection{Caso de estudio}

\begin{center}
\renewcommand{\arraystretch}{1.15}
\begin{tabular}{@{}llll@{}}
\toprule
\textbf{Parámetro} & \textbf{Símbolo} & \textbf{Valor} & \textbf{Observación}\\
\midrule
Altura del cilindro      & $H$      & \SI{0.250}{\meter}   & botella de \SI{1.26}{\liter}\\
Radio interior           & $R$      & \SI{0.040}{\meter}   & $A=\SI{50.3}{\centi\meter\squared}$\\
Masa de la botella vacía & $m_b$    & \SI{0.040}{\kilo\gram} & PET\\
Longitud de la cuerda    & $L_c$    & \SI{0.500}{\meter}   & $s=L_c+H=\SI{0.750}{\meter}$\\
Diámetro del orificio    & $d_a$    & \SI{5.0}{\milli\meter} & $a=\SI{19.6}{\milli\meter\squared}$\\
Coeficiente de descarga  & $C_d$    & $0.60$               & orificio de pared delgada\\
Altura inicial de agua   & $h_0$    & \SI{0.250}{\meter}   & botella llena, $m_{w0}=\SI{1.257}{\kilo\gram}$\\
Amplitud inicial         & $\Theta_0$ & \ang{12}           & \\
\midrule
Tiempo de vaciado        & $t_v$    & \SI{96.3}{\second}   & ec.~\eqref{eq:tv}\\
Período inicial          & $T_0$    & \SI{1.597}{\second}  & $\simeq 60$ oscilaciones\\
Adiabaticidad            & $\eps_{\max}$ & $9.2\times10^{-4}$ & ec.~\eqref{eq:epsilon}\\
\bottomrule
\end{tabular}
\end{center}

La integración numérica de \eqref{eq:sistema} se hizo con un método Runge--Kutta adaptativo de
orden 8 con tolerancias $10^{-12}$; el integrador se validó reproduciendo la corrección de
amplitud finita \eqref{eq:Tfinita} con error relativo $<2\times10^{-5}$.

\subsection{Vaciado}

\begin{figure}[H]
\centering
\begin{tikzpicture}
\begin{axis}[width=7.3cm,xlabel={$t$ (s)},ylabel={$h$ (mm)},xmin=0,xmax=96.5,ymin=0,
  legend pos=north east, ylabel style={blue!50!black}, y tick label style={blue!50!black},
  axis y line*=left, axis x line*=bottom]
\addplot[azulosc] table[x index=0,y index=1] {curvas.dat};
\label{p:h}
\end{axis}
\begin{axis}[width=7.3cm,xmin=0,xmax=96.5,ymin=0,ymax=1.35,axis y line*=right,axis x line=none,
  ylabel={$m_w$ (kg)},ylabel style={rojo},y tick label style={rojo},grid=none]
\addplot[rojo,dashed] table[x index=0,y index=2] {curvas.dat};
\end{axis}
\end{tikzpicture}
\hfill
\begin{tikzpicture}
\begin{axis}[width=7.3cm,xlabel={$t$ (s)},ylabel={$\gamma=\Gamma/I_O$ (s$^{-1}$)},
   xmin=0,xmax=96.5,ymin=0,scaled y ticks=false,
   yticklabel style={/pgf/number format/fixed,/pgf/number format/precision=3}]
\addplot[verde] table[x index=0,y index=8] {curvas.dat};
\end{axis}
\end{tikzpicture}
\caption{Izquierda: ley de vaciado de Torricelli \eqref{eq:ht}, cuadrática en $t$ (azul, altura;
rojo discontinuo, masa de agua). Derecha: tasa de frenado instantánea $\gamma=\Gamma/I_O$.
El frenado es \emph{máximo al principio}, cuando coinciden el gasto máximo y el brazo máximo entre
superficie libre y orificio, y se anula al agotarse el agua. En el último instante $\gamma$ cruza a
valores negativos (imperceptibles en la escala): es la inversión de signo de $\Gamma$ para
$h<h_c=R^2/8s$ anunciada tras la ec.~\eqref{eq:Gamma}.}
\label{fig:vaciado}
\end{figure}

\subsection{Evolución del período}

\begin{figure}[H]
\centering
\begin{tikzpicture}
\begin{axis}[xlabel={$t$ (s)},ylabel={$L_{\rm eq}$ (m)},xmin=0,xmax=96.5,
             ymin=0.62,ymax=0.73,legend pos=south west]
\addplot[azulosc] table[x index=0,y index=3] {curvas.dat};
\addplot[rojo,only marks,mark=*,mark size=1.9pt] coordinates {(61.14,0.7172)};
\node[rojo,anchor=north west,font=\scriptsize] at (axis cs:63,0.7150) {$s-h^\ast=L_{\rm eq}$};
\addplot[gris,dashed,forget plot] coordinates {(0,0.63399)(96.5,0.63399)};
\legend{$L_{\rm eq}(t)$,máximo}
\end{axis}
\end{tikzpicture}
\hfill
\begin{tikzpicture}
\begin{axis}[xlabel={$t$ (s)},ylabel={$T$ (s)},xmin=0,xmax=96.5,ymin=1.585,ymax=1.715,
             legend pos=south west,legend cell align=left]
\addplot[azulosc] table[x index=0,y index=4] {curvas.dat};
\addplot[rojo,only marks,mark=o,mark size=1.3pt] table[x index=0,y index=1] {ciclos.dat};
\legend{$2\pi\sqrt{L_{\rm eq}/g}$,{numérico, ciclo a ciclo}}
\end{axis}
\end{tikzpicture}
\caption{Izquierda: longitud equivalente $L_{\rm eq}=I_O/S$. Crece un $13\%$, alcanza su máximo en
$t^\ast=\SI{61.1}{\second}$ (cuando la superficie libre cruza el centro de oscilación, con
$h^\ast=\SI{33.4}{\milli\meter}$) y cae bruscamente al valor de la botella vacía en el último
tramo. La línea gris marca $L_{\rm eq}$ de la botella llena, prácticamente igual a la de la
botella vacía. Derecha: período. Los círculos son el período medido ciclo a ciclo en la
integración numérica completa; están sistemáticamente por encima de la curva lineal por la
corrección de amplitud finita \eqref{eq:Tfinita}, que los reproduce con error rms de
\SI{0.02}{\milli\second}.}
\label{fig:periodo}
\end{figure}

\begin{center}
\renewcommand{\arraystretch}{1.15}
\begin{tabular}{@{}lcccc@{}}
\toprule
 & $h$ (mm) & $L_{\rm eq}$ (m) & $T$ (s) & $s-h$ (m)\\
\midrule
Botella llena ($t=0$)              & 250.0 & 0.6340 & 1.5973 & 0.500\\
Mitad ($t=\SI{28.2}{\second}$)     & 125.0 & 0.6869 & 1.6627 & 0.625\\
\textbf{Máximo} ($t^\ast=\SI{61.1}{\second}$) & \textbf{33.4} & \textbf{0.7172} & \textbf{1.6989} & \textbf{0.717}\\
Botella vacía ($t=t_v$)            & 0.0   & 0.6346 & 1.5981 & 0.750\\
\bottomrule
\end{tabular}
\end{center}

\noindent
En la fila del máximo se aprecia la condición \eqref{eq:extremo}: $s-h^\ast = \SI{0.717}{\meter} =
L_{\rm eq}$, con coincidencia en las tres primeras cifras.

\begin{resultado}[Respuesta sobre el período]
El período \textbf{crece} un $6.4\%$, alcanza un máximo cuando queda sólo un $13\%$ del agua
(a los $0.63\,t_v$), y \textbf{decrece} después hasta recuperar, en el instante en que la botella
queda vacía, prácticamente el mismo valor que tenía llena ($\Delta T/T = 5\times10^{-4}$).
La curva $T(t)$ es asimétrica: sube lentamente durante casi dos tercios del vaciado y baja de
golpe en el último tercio.
\end{resultado}

\subsection{Evolución de la amplitud}

\begin{figure}[H]
\centering
\begin{tikzpicture}
\begin{axis}[xlabel={$t$ (s)},ylabel={amplitud $\Theta$ (grados)},xmin=0,xmax=96.5,
  ymin=1,ymax=400,ymode=log,legend pos=north west,legend cell align=left,width=10.2cm,height=6.1cm,
  ytick={1,2,5,10,20,50,100,200},yticklabels={1,2,5,10,20,50,100,200}]
\addplot[naranja,dashed] table[x index=0,y index=6] {curvas.dat};
\addplot[verde,dashed] table[x index=0,y index=7] {curvas.dat};
\addplot[gris,dotted,line width=1.2pt] table[x index=0,y index=1] {env_naive.dat};
\addplot[azulosc,line width=1.2pt] table[x index=0,y index=5] {curvas.dat};
\addplot[rojo,only marks,mark=o,mark size=1.2pt] table[x index=0,y index=1] {env_completo.dat};
\legend{(A) sólo bombeo adiabático,(B) sólo frenado del chorro,
        modelo ingenuo \eqref{eq:ELnaive},WKB completo \eqref{eq:amplitud},numérico}
\end{axis}
\end{tikzpicture}
\caption{Amplitud en escala logarítmica. Los dos mecanismos por separado son enormes y de signo
opuesto: el bombeo adiabático (A) llevaría la amplitud de \ang{12} a \ang{68}; el frenado del
chorro (B) la llevaría a \ang{1.5}. Su producto, la predicción WKB \eqref{eq:amplitud}
(azul), coincide con la envolvente de la integración numérica completa (círculos rojos) hasta la
tercera cifra: $\ang{8.64}$ frente a $\ang{8.64}$. El modelo lagrangiano ingenuo (gris punteado)
predice $\ang{70.8}$ y queda descartado de inmediato por cualquier observación.}
\label{fig:amplitud}
\end{figure}

\begin{figure}[H]
\centering
\begin{tikzpicture}
\begin{axis}[width=7.35cm,height=4.9cm,xlabel={$t$ (s)},ylabel={$\theta$ (grados)},xmin=0,xmax=8,ymin=-13,ymax=13,
             title={\footnotesize primeros \SI{8}{\second} (botella llena)},title style={font=\footnotesize}]
\addplot[azulosc,line width=0.6pt] table[x index=0,y index=1] {trazaA.dat};
\end{axis}
\end{tikzpicture}
\hfill
\begin{tikzpicture}
\begin{axis}[width=7.35cm,height=4.9cm,xlabel={$t$ (s)},ylabel={$\theta$ (grados)},xmin=88,xmax=96,ymin=-13,ymax=13,
             title={\footnotesize últimos \SI{8}{\second} (casi vacía)},title style={font=\footnotesize}]
\addplot[rojo,line width=0.6pt] table[x index=0,y index=1] {trazaB.dat};
\end{axis}
\end{tikzpicture}
\caption{Señal $\theta(t)$ al principio y al final de la maniobra, en la misma escala. Se aprecian
simultáneamente la reducción de amplitud (de \ang{12} a \ang{8.6}) y el hecho de que el período
al final ha vuelto a ser casi igual al del principio.}
\label{fig:traza}
\end{figure}

\begin{resultado}[Respuesta sobre la amplitud]
La amplitud \textbf{decrece de forma monótona}, de \ang{12} a \ang{8.64}, es decir, un factor
$0.720$. No hay ninguna fase de crecimiento porque el frenado del chorro domina sobre el bombeo
adiabático en todo instante. La ley completa es \eqref{eq:amplitud}, y el factor total,
$e^{-\mathcal{K}/2}$ con $\mathcal K = 0.658$, es independiente del diámetro del orificio.
\end{resultado}

\medskip
\noindent
Es importante subrayar que este amortiguamiento \textbf{no es disipación viscosa}: no hay fricción
en el modelo. Es un decaimiento puramente inercial, debido a que el sistema exporta momento
angular. Su firma experimental lo distingue del rozamiento con el aire: el rozamiento produce
decaimiento exponencial con tasa aproximadamente constante y \emph{continúa} después de que la
botella se haya vaciado, mientras que $\gamma(t)$ del modelo (Fig.~\ref{fig:vaciado}, derecha) se
anula exactamente en $t=t_v$.

%=======================================================================
\section{Regímenes no adiabáticos y efectos paramétricos}
%=======================================================================

\subsection{Cuando $t_v$ se hace comparable al período}
\label{sec:rapido}

La validez de \eqref{eq:wkb} exige $\eps=|\dot\omega|/\omega^2\ll1$, lo que equivale a
$t_v\gg T$. ¿Qué ocurre cuando el orificio es grande? Dos límites exactos acotan el problema.

\paragraph{Límite súbito.} Si $t_v\ll T$ la gravedad no tiene tiempo de actuar y
\eqref{eq:balance} se reduce a $\dfrac{\dd}{\dd t}\ln\dot\theta = -\Gamma/I_O$, de modo que
$\dot\theta$ se multiplica por $e^{-\mathcal K}$ mientras $\theta$ permanece congelado. El
resultado depende drásticamente de la \emph{fase} en que ocurra el vaciado:
\begin{equation}
\frac{\Theta_f}{\Theta_0}\Big|_{\rm súbito} =
\begin{cases}
1, & \text{si se vacía en un extremo } (\dot\theta=0),\\[2mm]
\dfrac{\omega_0}{\omega_f}\,e^{-\mathcal K}, & \text{si se vacía al pasar por el centro }(\theta=0).
\end{cases}
\end{equation}

\paragraph{Relación con el límite adiabático.} El resultado adiabático \eqref{eq:amplitud} es
\emph{exactamente la media geométrica} de los dos valores súbitos:
\begin{equation}
\sqrt{\frac{\omega_0}{\omega_f}}\;e^{-\mathcal K/2}
=\sqrt{1\cdot\frac{\omega_0}{\omega_f}e^{-\mathcal K}} ,
\end{equation}
como debe ser, puesto que el régimen lento promedia sobre todas las fases.

\begin{figure}[H]
\centering
\begin{tikzpicture}
\begin{axis}[xlabel={$t_v/T$},ylabel={amplitud final (grados)},xmode=log,
  xmin=0.012,xmax=70,ymin=5.5,ymax=12.8,legend pos=south east,legend cell align=left,
  width=11.5cm,height=6.4cm]
\addplot[rojo,mark=*,mark size=1.1pt] table[x index=0,y index=1] {subito.dat};
\addplot[azulosc,mark=square*,mark size=1.1pt] table[x index=0,y index=2] {subito.dat};
\addplot[verde,mark=triangle*,mark size=1.3pt] table[x index=0,y index=3] {subito.dat};
\addplot[black,dashed,forget plot] coordinates {(0.012,8.638)(70,8.638)};
\node[anchor=south west,font=\scriptsize] at (axis cs:0.02,8.7) {predicción adiabática $\ang{8.64}$};
\legend{fase inicial $0$ (extremo),fase $\pi/2$ (centro),fase $\pi/4$}
\end{axis}
\end{tikzpicture}
\caption{Amplitud final tras un vaciado completo, en función de la rapidez del vaciado y de la fase
de la oscilación en que se inicia. Para $t_v\gtrsim 3\,T$ las tres curvas colapsan sobre la
predicción adiabática (línea negra), \emph{independientemente de $t_v$}: es la universalidad de la
Sec.~\ref{sec:universalidad}. Para $t_v\lesssim T$ el resultado se abre en abanico entre los dos
límites súbitos, \ang{12} y \ang{6.23}, y pasa a depender de la fase inicial. Notablemente, la
predicción adiabática sigue siendo buena dentro de $\pm 3\%$ hasta $t_v\approx T$.}
\label{fig:subito}
\end{figure}

\begin{resultado}[Robustez del régimen adiabático]
La transición entre ambos regímenes ocurre en $t_v \sim T$, no antes. Para la botella real, con
$t_v/T \approx 60$, el tratamiento adiabático es excelente. Para observar el régimen súbito haría
falta un orificio de $\gtrsim\SI{30}{\milli\meter}$ de diámetro, con el que la hipótesis de flujo
cuasiestacionario y la propia ley de Torricelli dejarían de ser válidas.
\end{resultado}

\subsection{¿Puede la modulación del vaciado producir efectos paramétricos?}
\label{sec:parametrico}

De \eqref{eq:geffmod}, la oscilación modula la gravedad efectiva en el orificio a frecuencia
$2\omega$, con profundidad relativa
\begin{equation}
\delta_g = \Theta^2\left(\frac{s}{2L_{\rm eq}}+\frac14\right)\approx 0.84\,\Theta^2 .
\end{equation}
Como $u\propto\sqrt{g_{\rm eff}}$, el gasto se modula con profundidad $\delta_g/2$. Integrando esta
ondulación durante medio ciclo, el nivel adquiere un rizado
\begin{equation}
\frac{\delta h}{h} \sim \frac{\delta_g}{2}\,\frac{|\dot h|/h}{2\omega},
\end{equation}
y la profundidad de modulación paramétrica de la frecuencia es
\begin{equation}
\eps_p = \left|\frac{\partial\ln\omega^2}{\partial\ln h}\right|\frac{\delta h}{h}
\;\sim\;
\underbrace{\left|\frac{\partial\ln\omega^2}{\partial\ln h}\right|}_{\le\,0.16}
\cdot\;\frac{0.42\,\Theta^2}{\omega}\cdot\frac{|\dot h|}{h}.
\label{eq:epsp}
\end{equation}
La ecuación \eqref{eq:lineal} adquiere así la forma de una ecuación de Mathieu amortiguada,
$\ddot\theta+\gamma\dot\theta+\omega^2[1+\eps_p\cos 2\omega t]\theta=0$, cuya lengua de resonancia
principal exige
\begin{equation}
\boxed{\;\eps_p > \frac{2\gamma}{\omega}\;}
\qquad\text{(umbral de crecimiento paramétrico).}
\end{equation}

\paragraph{Evaluación numérica.} Para el caso de estudio,
$\eps_p^{\max}=8.7\times10^{-6}$ mientras que $2\gamma/\omega\approx 3.7\times10^{-3}$: el sistema
está \textbf{dos órdenes y medio de magnitud por debajo del umbral}, y eso antes de añadir el
rozamiento del aire, que aumenta aún más el umbral.

\begin{resultado}[No hay resonancia paramétrica en el régimen lento]
La modulación del vaciado sí genera una estructura paramétrica genuina a frecuencia $2\omega$,
pero su profundidad escala como
\[
\eps_p \propto \frac{\Theta^2}{\omega}\,\frac{|\dot h|}{h}\sim\frac{\Theta^2}{\omega\,t_v},
\]
mientras que el umbral escala como $\gamma/\omega$. Un vaciado lento suprime el numerador sin
suprimir proporcionalmente el umbral. El bombeo paramétrico sólo podría competir si
$t_v\sim T$ \emph{y} $\Theta$ fuese grande, pero ése es exactamente el régimen en el que la
descripción de Mathieu (que presupone muchos ciclos) deja de tener sentido.
\textbf{La respuesta es negativa} para una botella real: la modulación paramétrica es un efecto
detectable en principio, pero no observable en la práctica.
\end{resultado}

%=======================================================================
\section{Límites de validez y correcciones descartadas}
\label{sec:validez}
%=======================================================================

Se cuantifican aquí las hipótesis de la Sec.~\ref{sec:ideal}, para justificar que las que se han
despreciado son realmente despreciables y para señalar dónde falla el modelo.

\subsection{Chapoteo (\emph{sloshing}): ¿es lícito congelar el agua?}

El primer modo antisimétrico de chapoteo en un cilindro de radio $R$ tiene frecuencia
\begin{equation}
\omega_{\rm sl}^2 = \frac{g\,\epsilon_{11}}{R}\tanh\!\left(\frac{\epsilon_{11}h}{R}\right),
\qquad \epsilon_{11}=1.8412 .
\end{equation}
Para $R=\SI{4}{\centi\meter}$ y columna alta, $\omega_{\rm sl}=\SI{21.2}{\radian\per\second}$
frente a $\omega=\SI{3.9}{\radian\per\second}$: una razón de $5.4$. El chapoteo está muy fuera de
resonancia y responde cuasiestáticamente.

\begin{center}
\begin{tabular}{@{}rccc@{}}
\toprule
$h$ (mm) & $\omega_{\rm sl}/\omega$ & $\Delta S_{\rm sup}/S$ (\%) & \\
\midrule
250 & 5.40 & 0.25 & régimen seguro\\
100 & 5.66 & 0.53 & régimen seguro\\
 20 & 4.89 & 2.02 & correcciones apreciables\\
  5 & 2.67 & 4.59 & modelo degradado\\
  1 & 1.17 & 6.99 & \textbf{modelo inválido}\\
\bottomrule
\end{tabular}
\end{center}

\noindent
La resonancia exacta $\omega_{\rm sl}=\omega$ se alcanza en $h=\SI{0.73}{\milli\meter}$, es decir,
con \SI{3.7}{\gram} de agua residual ($0.3\%$ del total). \textbf{La hipótesis de agua congelada es
excelente durante el $99.7\%$ del vaciado y falla sólo en el último instante}, donde de todos modos
la inercia del agua es irrelevante frente a la de la botella.

\subsection{Efecto de superficie libre}

Si la superficie libre se inclina un ángulo $\alpha$ respecto de la sección recta, el centro de
masa del agua se desplaza lateralmente $\delta x = (R^2/4h)\tan\alpha$, lo que reduce el momento
restaurador en
\begin{equation}
\Delta S_{\rm sup} \sim \rho\,\frac{\pi R^4}{4} = \SI{2.0e-3}{\kilo\gram\meter},
\end{equation}
el clásico ``efecto de superficie libre'' de la arquitectura naval. En nuestro caso es un
sobredimensionamiento, porque para un péndulo la gravedad efectiva es casi paralela al eje
(un plomo colgado de una cuerda cuelga a lo largo de la cuerda), de modo que $\alpha=O(\theta)\cdot
|1-r/L_{\rm eq}|\ll\theta$. Los valores de la tabla anterior son cotas superiores. Conclusión: por
debajo del $1\%$ mientras $h\gtrsim\SI{5}{\centi\meter}$.

\subsection{Otras correcciones}

\begin{itemize}[leftmargin=*,itemsep=2pt]
\item \textbf{Empuje del chorro.} $F = \rho a u^2 = \SI{34.7}{\milli\newton}$ frente a un peso de
\SI{12.7}{\newton}: $0.27\%$. Además está dirigido a lo largo del eje, luego no produce torque
respecto de $O$ (ya incluido exactamente en el balance).
\item \textbf{Contracción de la vena / número de Reynolds.} $Re\approx 1.1\times10^4$: flujo
turbulento plenamente desarrollado, $C_d$ estable y bien tabulado.
\item \textbf{Flexibilidad de la cuerda (vínculo V2).} Es la hipótesis más frágil del modelo. Una
cuerda real permite un modo de cabeceo adicional; se mitiga suspendiendo la botella con una
\emph{horquilla} de dos hilos en V, que impone V2 mecánicamente.
\item \textbf{Rozamiento del aire.} No está en el modelo. Para una botella de $\SI{8}{\centi\meter}$
de diámetro a $\SI{0.5}{\meter\per\second}$ es comparable al frenado por vaciado, y debe medirse
por separado (Sec.~\ref{sec:experimento}).
\item \textbf{Tensión superficial y entrada de aire.} Para orificios de $\lesssim\SI{2}{\milli\meter}$
el chorro puede volverse intermitente (glugueo) y la ley de Torricelli deja de valer.
Usar $d_a\ge\SI{4}{\milli\meter}$ y un respiradero superior.
\end{itemize}

%=======================================================================
\section{Protocolo experimental y contraste con el modelo}
\label{sec:experimento}
%=======================================================================

El modelo no tiene parámetros libres ajustables salvo $C_d$ (que sólo fija la escala temporal).
Esto permite un contraste exigente.

\subsection{Montaje}

Botella cilíndrica de sección constante (evitar botellas con hombros cónicos, que romperían
$A=$ cte y por tanto \eqref{eq:ht}); suspensión en horquilla de dos hilos para imponer V2;
orificio limpio de $\SI{4}{}$--$\SI{6}{\milli\meter}$ en el centro del fondo, con la tapa
\emph{abierta} o perforada para garantizar presión atmosférica en la superficie libre; marcador
de alto contraste en la botella; grabación a $\ge\SI{120}{fps}$ con la cámara perpendicular al
plano de oscilación y a varios metros de distancia (para minimizar la paralaje).

\subsection{Mediciones y análisis}

\begin{enumerate}[label=\textbf{M\arabic*.},leftmargin=*,itemsep=3pt]
\item \textbf{Calibración de $C_d$ sin oscilación.} Con la botella en reposo, medir $h(t)$ y
ajustar $\sqrt{h}$ frente a $t$: la ley \eqref{eq:ht} predice una recta cuya pendiente da
$C_d$. Conviene hacerlo antes, para que el vaciado quede totalmente predicho.

\item \textbf{Calibración del rozamiento del aire.} Repetir el experimento con la \emph{botella
sellada} y la misma masa de agua, midiendo el decaimiento de amplitud. Esto da la tasa de
amortiguamiento parásita $\gamma_{\rm aire}$, que debe restarse antes de comparar con
$\gamma=\Gamma/I_O$. \emph{Es la medición de control más importante del experimento.}

\item \textbf{Período ciclo a ciclo.} Detectar cruces por cero ascendentes del ángulo y calcular
$T_k$ como diferencia entre cruces consecutivos. Comparar con \eqref{eq:Tfinita}, no con la
fórmula lineal: la corrección $\Theta^2/16$ es del mismo orden que el efecto buscado.

\item \textbf{Localización del máximo.} Ajustar una parábola a la región del máximo de $T_k$ para
extraer $t^\ast$, y verificar la predicción geométrica \eqref{eq:extremo} midiendo la altura
$h^\ast$ del agua en ese instante en el vídeo. \textbf{Es el test más específico del modelo}:
predice un número (la posición de la superficie libre) a partir de otro completamente distinto
(el instante del máximo del período).

\item \textbf{Test de universalidad.} Repetir con orificios de $\SI{4}{}$, $\SI{5}{}$ y
$\SI{8}{\milli\meter}$. La predicción es que $t_v$ escala como $1/a$ pero que
$\Theta(t_v)/\Theta_0$ es \textbf{el mismo en los tres casos}. Un resultado independiente del
área del orificio confirmaría \eqref{eq:K}; una dependencia sistemática delataría que el
amortiguamiento observado es en realidad viscoso.

\item \textbf{Discriminación entre modelos.} Con $\Theta_0=\ang{12}$, el modelo ingenuo predice
$\ang{71}$ y el modelo corregido $\ang{8.6}$. No hace falta precisión: basta mirar.
\end{enumerate}

\subsection{Fuentes de discrepancia esperadas}

Si la amplitud medida decae más deprisa que $\ang{8.6}$, el exceso debe atribuirse al aire
(medido en M2) y al pivote. Si el período medido crece \emph{menos} del $6\%$ predicho, sospechar
de una botella no cilíndrica (hombros) o de un $h_0$ inicial menor que $H$. Si el máximo del
período aparece antes de $0.6\,t_v$, comprobar la posición del centro de masa de la botella vacía,
que es el parámetro al que \eqref{eq:extremo} es más sensible.

%=======================================================================
\section{Conclusiones}
%=======================================================================

\begin{enumerate}[leftmargin=*,itemsep=4pt]

\item El sistema tiene \textbf{un grado de libertad}, $\theta$, tras imponer seis vínculos, de los
cuales tres (V4--V6) son \textbf{reónomos}. El lagrangiano
$\mathcal L=\tfrac12 I_O(t)\dot\theta^2+gS(t)\cos\theta$ depende explícitamente del tiempo a través
de dos funciones escalares: el momento estático $S(t)$ y el momento de inercia $I_O(t)$.

\item \textbf{No hay coordenadas cíclicas} en el problema plano, y \textbf{ninguna cantidad se
conserva exactamente}: $\dd E/\dd t = -\partial\mathcal L/\partial t \ne 0$. En la generalización
tridimensional, $\varphi$ sí es cíclica, pero la pérdida de materia impide concluir que $p_\varphi$
se conserve: una lección general sobre el alcance del teorema de Noether en sistemas abiertos.

\item El lagrangiano con masa dependiente del tiempo es \textbf{insuficiente}. El balance de
momento angular con flujo de masa añade el término $\dot m_{\rm out}s^2\dot\theta$ y conduce a
\[
I_O\ddot\theta+\Gamma\dot\theta+gS\sin\theta=0,\qquad
\Gamma=\rho A|\dot h|\Big[h(2s-h)-\tfrac{R^2}{4}\Big]>0 .
\]
La diferencia entre ambos modelos es el $700\%$ en la amplitud final.

\item Con la ley de Torricelli $h(t)=h_0(1-t/t_v)^2$, \textbf{el período no es monótono}:
crece hasta un máximo y luego decrece, terminando en el valor de la botella vacía. El máximo ocurre
en el instante preciso en que \textbf{la superficie libre cruza el centro de oscilación},
$s-h^\ast=L_{\rm eq}(h^\ast)$. Para la botella de estudio, $+6.4\%$ con el máximo en $0.63\,t_v$.

\item \textbf{La amplitud decrece de forma monótona}, según
\[
\frac{\Theta(t)}{\Theta_0}=\left[\frac{S_0I_{O,0}}{S\,I_O}\right]^{1/4}
\exp\left[-\frac{s^2}{2}\int_0^t\frac{\dot m_{\rm out}}{I_O}\dd t'\right],
\]
producto de un bombeo adiabático ($\times5.7$) y un frenado por evacuación de momento angular
($\times0.127$). El resultado neto es $\times 0.72$.

\item El \textbf{invariante adiabático} $J=\pi\Theta^2\sqrt{gSI_O}$ es la única cantidad
casi-conservada, y su violación está dada exactamente por
$\dd J/\dd t = -(\dot m_{\rm out}s^2/I_O)\,J$.

\item El factor total de cambio de amplitud es \textbf{independiente del área del orificio}:
sólo depende de la integral geométrica \eqref{eq:K}. Es la predicción más falsable del modelo.

\item El modelo de \textbf{masa puntual predice período y amplitud constantes}: el fenómeno
completo es un efecto de cuerpo rígido extenso.

\item Los \textbf{efectos paramétricos} existen (modulación del gasto a $2\omega$ vía la gravedad
efectiva) pero quedan dos órdenes y medio por debajo del umbral de Mathieu. La aproximación
adiabática es válida hasta $t_v\approx T$, mucho más allá de lo que cabría esperar.
\end{enumerate}

\appendix
%=======================================================================
\section{Deducción de las identidades \eqref{eq:identidades}}
%=======================================================================

La contribución del agua al momento de inercia respecto de $O$ es
\begin{equation}
J_w(h) = \rho A h\left[\frac{R^2}{4}+\frac{h^2}{12}+\Big(s-\frac h2\Big)^2\right]
= \rho A\left[\frac{R^2 h}{4} + s^2 h - s h^2 + \frac{h^3}{3}\right],
\end{equation}
donde se ha usado $h\big(\frac{h^2}{12}+s^2-sh+\frac{h^2}{4}\big)
= \frac{h^3}{3}+s^2h-sh^2$. Derivando,
\begin{equation}
\frac{\dd J_w}{\dd h} = \rho A\left[\frac{R^2}{4}+s^2-2sh+h^2\right]
= \rho A\left[(s-h)^2+\frac{R^2}{4}\right],
\end{equation}
que es exactamente el momento de inercia de un disco de espesor $\dd h$, radio $R$ y masa
$\rho A\,\dd h$ situado a distancia $s-h$ del pivote (Steiner más el momento propio
$\frac14 m R^2$ de un disco). Análogamente,
$S_w=\rho A(sh-h^2/2)\Rightarrow \dd S_w/\dd h = \rho A(s-h)$. $\blacksquare$

%=======================================================================
\section{Del invariante adiabático a la ley de amplitud}
%=======================================================================

Para el oscilador $\ddot\theta+\gamma\dot\theta+\omega^2\theta=0$ con $\gamma,\omega$ lentos, la
sustitución $\theta=\psi\exp(-\frac12\int\gamma)$ y la aproximación WKB dan
$\Theta=\Theta_0\sqrt{\omega_0/\omega}\,e^{-\frac12\int\gamma}$.
Con $\gamma=\Gamma/I_O$ y $\Gamma=\dot I_O+\dot m_{\rm out}s^2$,
\begin{equation}
\exp\left[-\frac12\int_0^t\frac{\dot I_O}{I_O}\right]=\left(\frac{I_{O,0}}{I_O}\right)^{1/2},
\end{equation}
de donde, usando $\omega=\sqrt{gS/I_O}$,
\begin{equation}
\frac{\Theta}{\Theta_0}
=\left(\frac{\omega_0}{\omega}\right)^{1/2}\left(\frac{I_{O,0}}{I_O}\right)^{1/2}
e^{-\frac{s^2}{2}\int\frac{\dot m_{\rm out}}{I_O}}
=\left(\frac{S_0 I_{O,0}}{S\,I_O}\right)^{1/4}e^{-\frac{s^2}{2}\int\frac{\dot m_{\rm out}}{I_O}} .
\end{equation}
Sustituyendo en $J=\pi\Theta^2\sqrt{gSI_O}$ se obtiene $J=J_0\exp[-s^2\!\int\!\dot m_{\rm out}/I_O]$
y, derivando, la ecuación \eqref{eq:dJdt}. $\blacksquare$

\medskip
\noindent\textbf{Verificación del límite conservativo.} Si $\dot m_{\rm out}\to0$ con $I_O$ aún
variable (por ejemplo, una botella que se deforma sin perder masa), $\Gamma\to\dot I_O$ y
$J=$ cte: se recupera la conservación exacta de la acción para el lagrangiano \eqref{eq:L}.
Es la comprobación de que la corrección de masa variable no estropea el caso hamiltoniano.

%=======================================================================
\section{Guion de implementación numérica}
%=======================================================================

\begin{enumerate}[leftmargin=*,itemsep=1pt]
\item Definir $h(t)=h_0\max(0,1-t/t_v)^2$ y $\dot h(t)=-2h_0\max(0,1-t/t_v)/t_v$.
\item Definir $S(h)$, $I_O(h)$, $\Gamma(h,\dot h)$ según \eqref{eq:sistema}.
\item Integrar $\dot y_1=y_2$, $\dot y_2=-[\Gamma y_2 + gS\sin y_1]/I_O$ con RK adaptativo
(tolerancias $\le10^{-10}$; el paso debe resolver el período, no $t_v$).
\item Extraer el período de los cruces por cero ascendentes con interpolación lineal, y la
amplitud de los máximos locales.
\item \textbf{Validación obligatoria:} con $\dot h\equiv0$ el integrador debe reproducir
$T=T_0(1+\Theta^2/16+11\Theta^4/3072)$ con error relativo $<10^{-4}$.
\item Comparar la envolvente numérica con \eqref{eq:amplitud} y la secuencia de períodos con
\eqref{eq:Tfinita}.
\end{enumerate}

\end{document}
